% This is samplepaper.tex, a sample chapter demonstrating the
% LLNCS macro package for Springer Computer Science proceedings;
% Version 2.20 of 2017/10/04
%
\documentclass[runningheads]{llncs}
%
\newcommand{\kashima}[1]{\textcolor{red}{Kashima: {#1}}}
\newcommand{\harada}[1]{\textcolor{blue}{(Harada: {#1})}}
\newcommand{\haradap}[1]{\textcolor{green}{(Harada: {#1})}}

% Use the postscript times font!
\usepackage{times}
\usepackage[utf8]{inputenc}
\renewcommand*\ttdefault{txtt}
\usepackage{soul}
\usepackage{url}
\usepackage[hidelinks]{hyperref}
\usepackage[utf8]{inputenc}
\usepackage[small]{caption}
\usepackage{graphicx}
\usepackage{amsmath}
\usepackage{amsfonts}
\usepackage{booktabs}
\usepackage[ruled]{algorithm2e} % For algorithms
\usepackage{color}
\renewcommand{\algorithmcfname}{ALGORITHM}
\usepackage{subcaption,siunitx,booktabs}
\usepackage{multirow}
\urlstyle{same}
\usepackage{bm}

\newcommand{\x}{\mathbf x}
\usepackage{graphicx}
% Used for displaying a sample figure. If possible, figure files should
% be included in EPS format.
%
% If you use the hyperref package, please uncomment the following line
% to display URLs in blue roman font according to Springer's eBook style:
% \renewcommand\UrlFont{\color{blue}\rmfamily}

\sloppy

\begin{document}
%
\title{Counterfactual Propagation for Semi-Supervised Individual Treatment Effect Estimation}
%
\titlerunning{Counterfactual Propagation for Semi-supervised ITE Estimation}
% If the paper title is too long for the running head, you can set
% an abbreviated paper title here
%
% \author{First Author\inst{1} \and
% Second Author\inst{2}}
% \author{Shonosuke Harada\inst{1} \and Hisashi Kashima\inst{1,2}}
\author{Shonosuke Harada \and Hisashi Kashima}
% \author{Anonymous}
%
% \authorrunning{F. Author et al.}
% First names are abbreviated in the running head.
% If there are more than two authors, 'et al.' is used.
%
% \institute{Organization, Address, Country \email{email@domain}\and
% Organization, Address, Country
% \email{email@domain}}
\institute{Kyoto University\\ \email{sh1108@ml.ist.i.kyoto-u.ac.jp}}
%\institute{Kyoto University \and RIKEN AIP\\ %\email{sh1108@ml.ist.i.kyoto-u.ac.jp}\\ {kashima@i.kyoto-u.ac.jp}}
% \institute{1 Kyoto University, \and 2 RIKEN AIP\\
% \email{sh1108@ml.ist.i.kyoto-u.ac.jp}\\ \email{kashima@i.kyoto-u.ac.jp}}
% \institute{sh1108@ml.ist.i.kyoto-u.ac.jp, kashima@i.kyoto-u.ac.jp}


% \author{First Author\inst{1}\orcidID{0000-1111-2222-3333} \and
% Second Author\inst{2,3}\orcidID{1111-2222-3333-4444} \and
% Third Author\inst{3}\orcidID{2222--3333-4444-5555}}
% \authorrunning{F. Author et al.}
% First names are abbreviated in the running head.
% If there are more than two authors, 'et al.' is used.

% \institute{Princeton University, Princeton NJ 08544, USA \and
% Springer Heidelberg, Tiergartenstr. 17, 69121 Heidelberg, Germany
% \email{lncs@springer.com}\\
% \url{http://www.springer.com/gp/computer-science/lncs} \and
% ABC Institute, Rupert-Karls-University Heidelberg, Heidelberg, Germany\\
% \email{\{abc,lncs\}@uni-heidelberg.de}}

%
\maketitle              % typeset the header of the contribution
%
\begin{abstract}
Individual treatment effect (ITE) represents the expected improvement in the outcome of taking a particular action to a particular target, and plays important roles in decision making in various domains.
However, its estimation problem is difficult because intervention studies to collect information regarding the applied treatments (i.e., actions) and their outcomes are often quite expensive in terms of time and monetary costs. 
In this study, we consider a semi-supervised ITE estimation problem that exploits more easily-available unlabeled instances to improve the performance of ITE estimation using small labeled data.
We combine two ideas from causal inference and semi-supervised learning, namely, matching and label propagation, respectively, to propose {\it counterfactual propagation,} which is the first semi-supervised ITE estimation method. 
Experiments using semi-real datasets demonstrate that the proposed method can successfully mitigate the data scarcity problem in ITE estimation.

%\keywords{Causal Inference  \and Treatment effect estimation \and Semi-supervised learning.}
\end{abstract}
\section{Introduction\label{sec:introduction}}
\begin{figure}[tb]
 \begin{minipage}{0.33\hsize}
  \begin{center}
   \includegraphics[width=\linewidth]{images/toy_data_ja.pdf}
   (a):~Data distribution
  \end{center}
%   \caption{IHDP}
 \end{minipage}
 \begin{minipage}{0.33\hsize}
  \begin{center}
   \includegraphics[width=\linewidth]{images/OLS_1.pdf}
   (b):~Linear regression
  \end{center}
%   \caption{News}
 \end{minipage}
  \begin{minipage}{0.33\hsize}
  \begin{center}
   \includegraphics[width=\linewidth]{images/MLP.pdf}
    (c):~Neural networks
  \end{center}
%   \caption{IHDP}
 \end{minipage}
 \\
  \begin{minipage}{0.33\hsize}
  \begin{center}
 \includegraphics[width=\linewidth]{images/knn.pdf}
    (d):~$k$NN
  \end{center}
%   \caption{IHDP}
 \end{minipage}
  \begin{minipage}{0.33\hsize}
  \begin{center}
   \includegraphics[width=\linewidth    ]{images/LP.pdf}
   (e):~\bf Proposed method
  \end{center}
%   \caption{IHDP}
 \end{minipage}
   \begin{minipage}{0.07\hsize}
  \begin{center}
 \vspace{-0.72cm}
   \includegraphics[width=\linewidth   ]{images/cbar.pdf}
  \end{center}
%   \caption{IHDP}
 \end{minipage}
 \caption{Illustrative example using (a) two-moon dataset. Each moons has a constant ITE either of $0$ and $1$.
 Only two labeled instances are available for each moon, denoted by yellow points, whose observed $(\text{treatment}, \text{outcome})$ pairs are $(0, 1), (1, 1), (1, 1), (0,0)$ from left to right. 
Figures (b), (c), and (d) show the ITE estimation error (PEHE) by the standard two-model approach using different base models suffered from the lack of labeled data.
The deeper-depth color indicates larger errors. The proposed semi-supervised method (e) successfully exploits the unlabeled data to estimate the correct ITEs.\label{fig:toy} }
\end{figure}


One of the important roles of predictive modeling is to support decision making related to taking particular actions in responses to situations.
The recent advances of in the machine learning technologies have significantly improved their predictive performance.
However, most predictive models are based on passive observations and do not aim to predict the  causal effects of actions that actively intervene in environments.
%
For example, advertisement companies are interested not only in their customers' behavior when an advertisement is presented, but also in the causal effect of the advertisement, in other words, the change it causes on their behavior.
%
There has been a growing interest in moving from this passive predictive modeling to more active causal modeling in various domains, such as education~\cite{hill2011bayesian}, advertisement~\cite{chan2010evaluating,li2016matching}, 
economic policy~\cite{lalonde1986evaluating},
and health care~\cite{glass2013causal}.

Taking an action toward a situation generally depends on the expected improvement in the outcome due to the action.
% Generally speaking, whether or not we should take a particular action to a target depends on how much the outcome is expected to improve by taking the action.
This is often called the {\it individual treatment effect~(ITE)}~\cite{rubin1974estimating} and is defined as the difference between the outcome of taking the action and that of {\it not} taking the action.
%
An intrinsic difficulty in ITE estimation is that ITE is defined as the difference between the factual and counterfactual outcomes~\cite{lewis1974causation,pearl2009causality,rubin1974estimating}; in other words, the outcome that we can actually observe is either of the one when we take an action or the one when we do not, and it is physically impossible to observe both.
%
To address the counterfactual predictive modeling from observational data, various techniques including matching~\cite{rubin1973matching}, inverse-propensity weighting~\cite{rosenbaum1983central}, instrumental variable methods~\cite{baiocchi2014instrumental}, and more modern deep learning-based approaches have been developed~\cite{shalit2017estimating,johansson2016learning}.
For example, in the matching method, matching pairs of instances with similar covariate values and different treatment assignments are determined.
The key idea is to consider the two instances in a matching pair as the counterfactual instance of each other so that we can estimate the ITE by comparing the pair.


Another difficulty in ITE estimation is data scarcity.
For ITE estimation, we need some labeled instances whose treatments (i.e., whether or not an action was taken on the instance) and their outcomes (depending on the treatments) as well as their covariates are given.
However, collecting such labeled instances can be quite costly in terms of time and money, or owing to other reasons, such as physical and ethical constraints~\cite{radlinski2005query,pombo2015machine}. 
Consequently, ITE estimation from scarcely labeled data is an essential requirement in many situations.

In the ordinary predictive modeling problem, a promising option to the scarcity of labeled data is semi-supervised learning that exploits unlabeled instances only with covariates because it is relatively easy to obtain such unlabeled data. 
% since such unlabeled data is relatively easy to obtain.
A typical solution is the graph-based label propagation method~\cite{zhu2002learning,belkin2006manifold,weston2012deep}, which makes predictions for unlabeled instances based on the assumption that instances with similar covariate values are likely to have a same label.

In this study, we consider a semi-supervised ITE estimation problem.
The proposed solution called {\it counterfactual propagation} is based on the resemblance between the matching method in causal inference and the graph-based semi-supervised learning method called label propagation.
We consider a weighted graph over both labeled instances with treatment outcomes and unlabeled instances with no outcomes, and estimate ITEs using the smoothness assumption of the outcomes and the ITEs.

The proposed idea is illustrated in Fig.~\ref{fig:toy}.
Fig.~\ref{fig:toy}(a)~describes the two-moon shaped data distribution. 
We consider a binary treatment and binary outcomes.
The blue points indicate the instances with a positive ITE~($=1$), where the outcome is $1$ if the treatment is $1$ and $0$ if the treatment is $0$. 
The red points indicate the instances with zero ITE~($=0$); their outcomes are always $1$ irrespective of the treatments.
We have only four labeled data instances shown as yellow points, whose observed $(\text{treatment}, \text{outcome})$ pairs are $(0, 1), (1, 1), (1, 1), (0,0)$ from left to right. 
Since the amount of labeled data is considerably limited,
supervised methods relying only on labeled data fail to estimate the ITEs.
Figures~\ref{fig:toy}(b), (c), (d) show the ITE estimation errors by the standard two-model approach using different base learners, which show poor performance. 
In contrast, the proposed approach exploits unlabeled data to find connections between the red points and those between the blue points to estimate the correct ITEs (Fig.~\ref{fig:toy}(e)).

\if0
\harada{
% 本研究の動機をより明確にするために, 人工データを用いた実験を行う. 我々の関心は, 利用可能な教師付きインスタンスが厳しく制限された状況下で, 教師付きインスタンスのみを用いた通常の学習ではどのような結果を招くのかという点にある. 図~\ref{fig:toy}において半教師付き学習問題で典型的なtwo moon型のデータ分布による実験の結果を示す. 図~\ref{fig:toy}(a)~は人工データの分布であり, 赤と青の点はそれぞれ正の介入効果がある~(ITE$=1$)~インスタンス及び介入効果がない~(ITE$=0$)~インスタンスを示す. 簡単の為, 正の介入効果を持つ全てのインスタンスに対する介入効果を$1$とする.  図~\ref{fig:toy}(b), (c), (d), (e)はそれぞれ線形回帰, ニューラルネットワーク, k近傍法~(KNN)~, 提案手法による実験結果を示す. ヒートマップは誤差の程度を示しており色が濃い程大きな誤差を生んでいることを示す. $3$つのモデルを単純な教師付き学習によって構築し介入効果を予測すると, それぞれのモデルは特定のインスタンスについては大きな誤差を示す. 一方, 半教師付き学習によって予測を行うと, ほとんど全てのインスタンスに対して一定の予測性能を示しており, 教師インスタンスが少数である場合でも一定頑健である. 我々はこのような現象が実世界でも生じうることを考慮し, 教師付きインスタンスの数が厳しく制限されている状況下でも頑健なモデルを構築することに動機づけされている.
We are particularly interested in examining what will happen if we just employ just a supervised model when the available labeled data are strictly limited. Fig.~\ref{fig:toy}(a)~describes the distribution of two moon dataset. Yellow points represent labeled data whose covariates, type of treatments, and outcomes have been observed. Blue and red points represent instances which have positive ITE~(ITE$=1$)~ and no ITE~(ITE$=0$), respectively. For simplicity, we set the all ITE to the same value, $1$. Fig.(b)-(e) represent experimental results using linear regression, neural networks, k nearest neighbors~(k-NN), and proposed method. If we train widely employed models in just supervised manner, we observed these models showed significantly poor performance for some instances~(Figure~\ref{fig:toy}(a), \ref{fig:toy}(b), \ref{fig:toy}(c), \ref{fig:toy}(d)). In contrast, our proposed method showed moderate performance for almost all instances. We consider similar phenomena could be possible in the real-world and are highly motivated to build a model which can maintain robustness under such a severe situation.
% Inspired by counterfactual problem, we are interested in what will happen
% under situations when available labeled instances are strictly limited  stated in Section~\ref{sec:introduction}. Figure~\ref{fig:toy}~describes experimental results using toy data based on two moon. Red and blue are instances which have no treatmen effect~(${\rm ITE}=0$)~and positive treatment effect~(${\rm ITE}=1$), respectively. For simplicity, we set the all positive treatment effects to $1$. If we train widely employed models in just supervised manner, we observe these models showed significantly poor performances for some instances~(Figure~\ref{fig:toy}a, \ref{fig:toy}b, \ref{fig:toy}c, \ref{fig:toy}d). In contrast, our proposed model showed moderate performance for almost all instances. We assume similar phenomena could be possible in the real-world and are highly motivated to build a model which can maintain robustness in such a severe situation.
}
\fi


We propose an efficient learning algorithm assuming the use of a neural network as the base model, and conduct experiments using semi-synthetic real-world datasets to demonstrate that the proposed method estimates the ITEs more accurately than baselines when the labeled instances are limited.


\if0
==========
Our goal is to estimate the treatment effect per instance, which is referred to as {\it individual treatment effect~(ITE)~} in this context~\cite{rubin1974estimating}. %However, we face a challenge in this problem.
%ITE is the difference between 
Treatment effect estimation problem has completely different aspects and challenges from a standard supervised learning. The ground truth in this problem is the difference between the treatment outcome and control outcome per instance. However, we can only access to factual outcomes and never observe outcomes of  actions which was not been actually conducted. This is called {\it counterfactual}~\cite{lewis1974causation,pearl2009causality,rubin1974estimating}.  Second, in practical, the treatment assignment often tend to be biased. That is specific instances who are considered to achieve desirable results often receive treatments and vice versa. Hence, standard supervised methods can not be straightforwardly applied in this problem. 
% Randomized controlled trials~(RCT)~is one of the approaches in which the distribution of treatment and control group is known,   

Moreover, in order to estimate treatment effect, we need some instances as labeled data whose treatments and outcomes have been observed as well as confounder variables. In some domains,  collecting such instances would be seriously time consuming and expensive~\cite{radlinski2005query,pombo2015machine}. Therefore, it is not hard to imagine that there would be cases when the number of available labeled data is strictly limited. Although labeled instances are limited, we can easily access to confounder variables of unlabeled data. For example, in health care, even though we do not have much labeled data, we can know several confounder variables such as age, sex, and blood type of patients who we would like to prescribe. In this way, the limitation of available labeled data and the existence of observable confounder variables are both realistic scenarios in many domains. Under such a situation, semi-supervised learning can be naturally introduced as a practical setting.




%We can not always observe the outcome of other treatments which is called {\it counterfactual} .

%{\it counterfactual problem}. Counterfactual problem means we can only observe one factual outcome per one unit.  

%Note that we make "no-hidden confounders" assumption, which assumes there no confounder variables affecting results. This assumption can be formalized using the strong ignorability condition:$(y_0,y_1),t\mid x$, and $0<p(t=1\mid x)<1$. 
 Semi-supervised learning is a very standard learning approach which makes use of both labeled data and unlabeled data. Semi-supervised learning have been quite effective especially when we are given limited number of labeled instances~\cite{bengio2007greedy,hinton2006fast}. 
 In many problems of semi-supervised learning methods, they try to leverage abundant unlabeled instances to improve predictive performance based on some assumed relationships between labeled and unlabeled instances. One of the assumption is that similar instances would have similar outcomes~\cite{zhu2002learning,belkin2006manifold,kipf2016semi}. 
 
 In treatment effect estimation problem, traditional matching based methods, which construct pairs of instances and predict counterfactual outcomes, have been developed on the similar assumption. Hence, solving treatment effect estimation problem over the framework of semi-supervised learning does not seem unnatural and undoubtedly important for the reasons as mentioned.  However, to the best of our knowledge, none of existing methods have not tried to consider the treatment effect estimation problem based on the semi-supervised learning problem setting. 

In this paper, we newly formalize the treatment effect estimation problem as a semi-supervised learning problem and aim to improved semi-supervised technique for mitigating counterfactual and limited labeled data problems  in a non-trivial manner. Our contributions in this work are summarized as follows:
\begin{itemize}
\item We introduce a new problem setting, semi-supervised treatment effect estimation, which exploits not only labeled instances but also unlabeled instances to improve model performance.
\item We propose a novel semi-supervised treatment effect estimation method based on the regularization which encourages the model to have similar outcomes for similar inputs. We extend this idea to treatment effect estimation problem.
\item On experiments using the semi-synthetic real-world datasets, we demonstrate our method works better especially in the case where the available labeled instances are limited.
\end{itemize}


\fi
%\section{Preliminaries\label{sec:problem}}

\section{Semi-supervised ITE estimation problem\label{sec:problem}}
We start with the problem setting of the semi-supervised treatment effect estimation problem. 
Suppose we have $N$ labeled instances and $M$ unlabeled instances. (We usually assume $N \ll M$.)
The set of labeled instances is denoted by $\{(\x_i, t_i, y^{t_i}_i)\}_{i=1}^{N}$, where $x_i \in \mathbb{R}^D$ is the covariates of the $i$-th instance, $t_i \in \{0,1\}$ is the treatment applied to instance $i$, and $y^{t_i}_i$ is its outcome.
Note that for each instance $i$, either  
$t_i=0$ or $t_i=1$ is realized; accordingly, either $y^{0}_i$ or $y^{1}_i$ is available. The unobserved outcome is called a counterfactual outcome.
%
The set of unlabeled instances is denoted by $\{(\x_i)\}_{i=N+1}^{N+M}$, where only the covariates are available.

Our goal is to estimate the ITE for each instance.
Following the Rubin-Neyman potential outcomes framework~\cite{rubin1974estimating,splawa1990application}, the ITE for instance $i$ is defined as   $\tau_i=y^1_{i}-y^0_{i}$
exploiting both the labeled and unlabeled sets. 
Note that $\tau_i$ is not known even for the labeled instances, and we want to estimate the ITEs for both the labeled and unlabeled instances.

%We focus on a binary treatment case for simplicity, but our proposed method  can be extended to multiple treatment setting. 

% \subsection{Motivation}
% Inspired by counterfactual problem, we are interested in what will happen
% under situations when available labeled instances are strictly limited  stated in Section~\ref{sec:introduction}. Figure~\ref{fig:toy}~describes experimental results using toy data based on two moon. Red and blue are instances which have no treatmen effect~(${\rm ITE}=0$)~and positive treatment effect~(${\rm ITE}=1$), respectively. For simplicity, we set the all positive treatment effects to $1$. If we train widely employed models in just supervised manner, we observe these models showed significantly poor performances for some instances~(Figure~\ref{fig:toy}a, \ref{fig:toy}b, \ref{fig:toy}c, \ref{fig:toy}d). In contrast, our proposed model showed moderate performance for almost all instances. We assume similar phenomena could be possible in the real-world and are highly motivated to build a model which can maintain robustness in such a severe situation.

%\subsection{Assumptions}
We make typical assumptions in ITE estimation in this study. i.e.,~({\rm i})~stable unit treatment value: the outcome of each instance is not affected by the treatment assigned to other instances;~({\rm ii})~unconfoundedness: the treatment assignment to an instance is independent of the outcome given covariates (confounder variables); ~({\rm iii})~overlap: each instance has a positive probability of treatment assignment.


\section{Proposed method\label{sec:proposed}}
We propose a novel ITE estimation method that utilizes both the labeled and unlabeled instances.
The proposed solution called {\it counterfactual propagation} is based on the resemblance between the matching method in causal inference and the graph-based semi-supervised learning method.
%We consider a weighted graph over both labeled instances with treatment outcomes and unlabeled instances with no outcomes, and estimate ITEs using the smoothness assumption of the outcomes and the ITEs.
%
% We make use of several semi-supervised learning techniques in order to estimate treatment effect even under the number of labeled instances are limited. 

\subsection{Matching}
Matching is a popular solution to address the counterfactual outcome problem.
Its key idea is to consider two similar instances as the counterfactual instance of each other so that we can estimate the causal effect by comparing the pair.
More concretely, we define the similarity $w_{ij}$ between two instances $i$ and $j$, as that defined between their covariates; for example, we can use the Gaussian kernel.
\begin{equation}\label{eq:GK}
    w_{ij} = \exp \left(- \frac{\| \x_i - \x_j \|^2}{\sigma^2} \right).
\end{equation}
The set of $(i,j)$ pairs with $w_{ij}$ being larger than a threshold and satisfying $t_i \neq t_j$ are found and compared as counterfactual pairs.
%
Note that owing to definition of the matching pair, the matching method only uses labeled data.

\subsection{Graph-based semi-supervised learning}\label{sec:gb}
Graph-based semi-supervised learning methods assume that the nearby instances in a graph are likely to have similar outputs. 
For a labeled dataset $\{(\x_i, y_i)\}_{i=1}^N$ and an unlabeled dataset $\{\x_i\}_{i=N+1}^{N+M}$, their loss functions for standard predictive modeling typically look like
\begin{equation}\label{eq:LPreg}
L(f)=\sum_{i=1}^{N}l(y_i, f(\x_i))+\lambda\sum_{i,j=1}^{N+M}w_{ij}\left(f(\x_i)-f(\x_j) \right)^2,
\end{equation}
where $f$ is a prediction model, $l$ is a loss function for the labeled instances, and $\lambda$ is a hyper-parameter.
The second term imposes ``smoothness" of the model output over the input space characterized by $w_{ij}$ that can be considered as the weighted adjacency matrix of a weighted graph; it can be seen the same as that used for matching~(\ref{eq:GK}).

The early examples of graph-based methods include label propagation~\cite{zhu2002learning} and manifold regularization~\cite{belkin2006manifold}.
More recently, deep neural networks have been used as the base model $f$~\cite{weston2012deep}.

% あまり本意ではないが, 文量が埋まらなさそうな気がしてきたので, TARNetの説明を入れて伸ばす
\subsection{Treatment effect estimation using neural networks}
% Recently, deep-learning based methods~\cite{shalit2017estimating,johansson2016learning} have been successfully applied in treatment effect estimation, we built our model based on deep learning framework~\cite{weston2012deep}.
%Recently, deep-learning based methods have been successfully applied in treatment effect estimation. 
We build our ITE estimation model based on the recent advances of deep-learning approaches for ITE estimation, specifically, the treatment-agnostic representation  network~(TARNet)~\cite{shalit2017estimating} that is a simple but quite effective model. 
%\kashima{Is the following sentence correct?}
% \harada{TARNet shares common parameters for both treatment instances and control instances, but employs different parameters in its prediction layers},
TARNet shares common parameters for both treatment instances and control instances to construct representations but employs different parameters in its prediction layer,
which is given as:
% \begin{equation}
    \begin{align}\label{eq:tarnet}
    f({\bf x}_i, t_i)=\left\{ \begin{array}{ll}
    \Theta_1^{\top}g\left(\Theta^{\top}x_i\right)&\ (t_i=1)\\
    \Theta_0^{\top}g\left(\Theta^{\top}x_i\right) &\ (t_i=0)\\
  \end{array}, \right.
    \end{align}
% \end{equation}
where $\Theta$ is the parameters in the representation learning layer and $\Theta_1, \Theta_0$ are those in the  prediction layers for treatment and controlled instances, respectively. 
The $g$ is a non-linear function such as ReLU. 
One of the advantages of TARNet is that joint representations learning and separate prediction functions for both treatments enable more flexible modeling. 


% Graph based regularization is a technique which has been applied in semi-supervised learning tasks. Utilizing a given or a generate graph based on nearest neighbors, graph based regularization encourage neighbor instances to have similar labels or outcomes based on the assumption that nearby instances in the graph should have similar results. Label Propagation~\cite{zhu2002learning} 

\subsection{Counterfactual propagation}
% As we have already seen, the

It is evident that the matching method relies only on labeled data, while the graph-based semi-supervised learning method does not address ITE estimation; however, 
they are quite similar because they  both use instance similarity to interpolate the  factual/counterfactual outcomes or model predictions as mentioned in Section \ref{sec:gb}.
Our idea is to combine the two methods to propagate the outcomes and ITEs over the matching graph assuming that similar instances would have similar outcomes.
% we combine the matching method with graph-based regularization.

Our objective function consists of three terms, $L_\text{s}, L_\text{o}, L_\text{e}$, given as
\begin{equation}\label{eq:proposed}
L(f) = L_\text{s}(f) +\lambda_\text{o}L_\text{o}(f) +\lambda_\text{e}L_\text{e}(f),
\end{equation}
where $\lambda_o$ and $\lambda_e$ are the regularization hyper-parameters. 
%the three main terms are~$({\rm i})$~a supervised outcome regression loss,~$({\rm ii})$~a neighbor outcome regularizer,~$({\rm iii})$~a neighbor treatment effect regularizer. 
% \harada{
We employ TARNet~\cite{shalit2017estimating} as the outcome prediction model $f(\x, t)$.
% \kashima{What do the next two sentences mean?? (Nihon-go de OK) $\rightarrow$ }
% One of the advantages of using TARNet is that joint learning of representations of each treatment makes the model more flexible in both obtaining representation and better predictive performance. 
% When it comes to propagating outcomes and ITEs, handling both of them in the single model lets the model propagate easier and more effectively.}\haradap{Ittan Nihongo de kaitemimasita (file sansho)}
% \harada{treatmentのインスタンスもcontrolのインスタンスも同一のモデルでrepresentation learningを行えるため, より柔軟な表現が学習でき, それぞれで別々にrepresentationおよびoutcomeを予測するよりも多くのデータを扱うことができ, 良い結果になると期待される. また類似インスタンスのoutcomeを考慮するラベル伝搬では, 一層期待される. } とか書いてみましたが、消した方がよい気がするので消します。
The first term in the objective function (\ref{eq:proposed})  is a standard loss function for supervised outcome estimation; we specifically employ the squared loss function as
\begin{equation}
% \vspace{-3mm}
\label{eq:supervised}
L_s(f)=\sum_{i=1}^{N}(y^{t_i}_i-f(\x_i, t_i))^{2}.
\end{equation}
% We employ the treatment-agnostic representation network~(TARNet)~as the base neural network model $f$~\cite{shalit2017estimating} that shares common parameters for the treatment and control instances but employs different parameters for the predictive layers. 
%The difference from the standard supervised loss is 
Note that it relies only on the observed outcomes of the treatments that are observed in the data denoted by $t_i$.
%Note that if we set $\lambda_o=\lambda_e=0$, our model reduces to TARNet~\cite{shalit2017estimating}. 

The second term $L_\text{o}$ is the outcome propagation term:
\begin{equation}
\label{eq:outcome}
\begin{split}
L_\text{o}(f)=\sum_{t}\sum_{i,j=1}^{N+M}w_{ij}((f(\x_i, t)-f(\x_j, t))^{2}.
\end{split}
\end{equation}
Similar to the regularization term (\ref{eq:LPreg}) in the graph-based semi-supervised learning, this term encourages the model to output similar outcomes for similar instances by penalizing the difference between their outcomes. 
This regularization term allows the model to propagate outcomes over a matching graph.
If two nearby instances have different treatments, they interpolate the counterfactual outcome of each other, which compares the factual and (interpolated) counterfactual outcomes to estimate the ITE.
%In this problem, even labeled instances have missing outcomes~(counterfactual)~as well as unlabeled instances. Thereby, differently from standard graph based regularization, by giving regularization for both treatment and control cases, we expect to build a better predictive model for treatment effect as well as outcome.
The key assumption behind this term is the smoothness of outcomes  for each treatment over the covariate space. 
While $w_{ij}$ indicates the adjacency between nodes $i$ and $j$  in the graph-based regularization, it can be considered as a matching between the instances $i$ and $j$ in the treatment effect estimation problem.
Even though traditional matching methods have only rely on labeled instances, we combine matching with graph-based regularization which also utilizes unlabeled instances. 
This regularization enables us to propagate the outcomes for each treatment over the matching graph and mitigate the counterfactual problem. 


% Utilizing scheme of related work~\cite{weston2012deep}, we introduce the following loss function:
% \begin{equation}\label{eq:outcome2}
% L_o=\sum_{t}\sum_{i,j=1}^{N+M}W_{i,j}((g(x_i, t)-g(x_j, t))^{2},
% \end{equation}
% where $g$ is a function which shares the first $m$ layer with supervised network and has new weights.

% We believe this term have a large difference from standard graph laplacian regularization losses since standard setting do not assume several outcomes for each instance if they are exposed to the same situation.
The third term $L_\text{e}$ is the ITE propagation term defined as
\begin{align}\label{eq:treatment}
L_\text{e}(f)=  \sum_{i,j=1}^{N+M}w_{ij} (\hat{\tau}_i - \hat{\tau}_j)^{2},
\end{align}
where $\hat{\tau}_i$ is the ITE estimate for instance $i$:
\[
\hat{\tau}_i = f(\x_i, 1)-f(\x_i, 0).
\]
\if0
\begin{align}\label{eq:treatment}
& L_\text{e}(f)= \nonumber \\
&~~ \sum_{i,j=1}^{N+M}w_{ij} ((\underbrace{f(x_i, 1)-f(x_i, 0)}_\text{ITE for instance $i$}) - (\underbrace{f(x_j, 1)-f(x_j, 0)}_\text{ITE for instance $j$}))^{2},
\end{align}
\fi
$L_\text{e}(f)$ imposes the smoothness of the ITE values in addition to that of the outcomes  imposed by outcome propagation (\ref{eq:outcome}).
In comparison to the standard supervised learning problems, where the goal is to predict the outcomes, as stated in Section~\ref{sec:problem}, our objective is to predict the ITEs.
This term encourages the model to output similar ITEs for similar instances. 
%Hence, based on the assumption that similar instances would have similar outcomes under the same treatment, we can consider the smoothness of treatment effect as well as the outcome. 
%The treatment effect for an instance $i$ predicted by $f$ is given as $f(x_i, 1)-f(x_i, 0)$. Following the smoothness assumption on treatment effect, we assume similar instances would have similar treatment effects with regard to each treatment. That is the difference of treatment effect between an instance $i$ and $j$, $(f(x_i, 1)-f(x_i, 0))-(f(x_j, 1)-f(x_j, 0))$, should be smaller.  
We expect that the outcome propagation and ITE propagation terms are beneficial especially when the available labeled instances are limited while there is an abundance of unlabeled  instances, similar to semi-supervised learning.


\if0
More specifically, because $f(x_i, 1)-f(x_i, 0)$ indicates predicted ITE for an instance $i$, ${\rm\widehat{ITE}}_i$, this regularization term can be interpreted as:
\begin{equation}
\label{eq:treatment}
L_e=\sum_{i,j=1}^{N+M}W_{ij} ({\rm\widehat{ ITE}}_i-{\rm \widehat{ITE}}_j) ^{2}.
\end{equation}
%We believe this regularization is very identical in this treatment effect estimation problem where the goal is not only to predict outcomes but effects.
In this paper, since we only focus on the binary treatment setting for simplicity, we introduce the treatment effect regularization using binary intervention. However our method can be easily extended to multiple intervention setting. 
\fi




\if0
\subsection{Supervised outcome regression loss}

The first term is the standard supervised loss term. 
\begin{equation}
% \vspace{-3mm}
\label{eq:supervised}
L_s=\sum_{i=1}^{N}(y^{t_i}_i-f(x_i, t_i))^{2},
\end{equation}
where $f$ is a predictive base model such as neural networks. We employ Treatment-Agnostic Representation Networks~(TARNet)~as a base neural network model~\cite{shalit2017estimating} which shares parameters for treatment and control instances while employs different parameters for predictive layers. The difference in this term from the standard supervised loss is that we only access to observed outcome in each treatment which is denoted by $t_i$ for each observed instance. 


\subsection{Outcome regularization}
The second term is the outcome regularization term. This encourages the model to output similar outcomes for similar instances by penalizing when similar instances have different outcomes. This regularization term allows the model to propagate outcomes over a matching graph. In this problem, even labeled instances have missing outcomes~(counterfactual)~as well as unlabeled instances. Thereby, differently from standard graph based regularization, by giving regularization for both treatment and control cases, we expect to build a better predictive model for treatment effect as well as outcome.
\begin{equation}
\label{eq:outcome}
\begin{split}
L_o=\sum_{t}\sum_{i,j=1}^{N+M}W_{ij}((f(x_i, t)-f(x_j, t))^{2}.
\end{split}
\end{equation}
The key assumption behind this term is the smoothness of outcomes in covariate space for each treatment. While in the graph-based regularization, $w_{ij}$ indicates the adjacency between node $i$ and $j$, in treatment effect estimation problem $w_{ij}$ can be considered as a matching between an instance $i$ and $j$. Though traditional matching methods have been only rely on labeled instances, we combine matching with graph-based regularization which utilizes unlabeled instances. This regularization enables us to propagate outcomes for each treatment over the matching graph and mitigate the counterfactual problem. We expect that the this model benefits from the propagation of outcomes especially in available instances are severely limited while observable instances are a lot.

% Utilizing scheme of related work~\cite{weston2012deep}, we introduce the following loss function:
% \begin{equation}\label{eq:outcome2}
% L_o=\sum_{t}\sum_{i,j=1}^{N+M}W_{i,j}((g(x_i, t)-g(x_j, t))^{2},
% \end{equation}
% where $g$ is a function which shares the first $m$ layer with supervised network and has new weights.

\subsection{Treatment effect regularization}
% We believe this term have a large difference from standard graph laplacian regularization losses since standard setting do not assume several outcomes for each instance if they are exposed to the same situation.
The third term is the treatment effect regularization term. This encourages the model to output similar treatment effect for similar instances. In standard supervised learning problems, their target is just to predict outcomes such as real values or labels for each instance. However, as stated in Section~\ref{sec:problem}, our goal is not only predict outcomes but also treatment effects for each instance. Hence, based on the assumption that similar instances would have similar outcomes under the same treatment, we can consider the smoothness of treatment effect as well as the outcome. The treatment effect for an instance $i$ predicted by $f$ is given as $f(x_i, 1)-f(x_i, 0)$. Following the smoothness assumption on treatment effect, we assume similar instances would have similar treatment effects with regard to each treatment. That is the difference of treatment effect between an instance $i$ and $j$, $(f(x_i, 1)-f(x_i, 0))-(f(x_j, 1)-f(x_j, 0))$, should be smaller.  In the same way as outcome regularization term, we introduce the following loss function:
\begin{equation}\label{eq:treatment}
L_e=\sum_{i,j=1}^{N+M}W_{ij} ((f(x_i, 1)-f(x_i, 0)) - (f(x_j, 1)-f(x_j, 0))) ^{2}.
\end{equation}
\noindent
More specifically, because $f(x_i, 1)-f(x_i, 0)$ indicates predicted ITE for an instance $i$, ${\rm\hat{ITE}}_i$, this regularization term can be interpreted as:
\begin{equation}
\label{eq:treatment}
L_e=\sum_{i,j=1}^{N+M}W_{ij} ({\rm\hat{ ITE}}_i-{\rm \hat{ITE}}_j) ^{2}.
\end{equation}
We believe this regularization is very identical in this treatment effect estimation problem where the goal is not only to predict outcomes but effects.
In this paper, since we only focus on the binary treatment setting for simplicity, we introduce the treatment effect regularization using binary intervention. However our method can be easily extended to multiple intervention setting. 



\subsection{Semi-supervised treatment effect estimation }
Finally the objective function of our proposed method is introduced by combining the regularization terms above:
\footnotesize
\begin{equation}\label{eq:proposed}
\begin{split}
L&=\overbrace{\sum_{i=1}^{N}(y^{t_i}_i-f(x_i, t_i))^{2}}^{\text{Supervised loss}}\\&\quad+\lambda_o\overbrace{\sum_{t}\sum_{i,j=1}^{N+M}W_{ij}((f(x_i, t)-f(x_j, t))^{2}}^{\text{Outcome regularization}}\\ &\quad+\lambda_e\overbrace{\sum_{i,j=1}^{N+M}W_{ij}((f(x_i, 1)-f(x_i, 0)) - (f(x_j, 1)-f(x_j, 0)))^2}^{\text{Treatment effect regularization}}\\
&=L_s+\lambda_{o}L_o+\lambda_{e}L_e,
\end{split}
\end{equation}\normalsize
where $\lambda_o, \lambda_e$ are the regularization parameters which control the strengths of regularization terms. 
Note that if we set $\lambda_o=\lambda_e=0$, our model reduces to TARNet~\cite{shalit2017estimating}. 
\fi

\subsection{Estimation algorithm}
As mentioned earlier, we assume the use of neural networks as the specific choice of the outcome prediction model $f$ based on the recent successes of deep neural networks in causal inference. 
For computational efficiency, we apply a sampling approach to optimizing Eq.~(\ref{eq:proposed}).
Following the existing method~\cite{weston2012deep}, we employ the Adam optimizer~\cite{kingma2014adam},~which is based on stochastic gradient descent to train the model  in a mini-batch manner. 

Algorithm~\ref{alg:train}~describes the procedure of model training, which iterates two steps until convergence. 
% In the first step, we sample a mini-batch consisting of $b_1$ labeled instances to approximate the supervised loss (\ref{eq:supervised}) and update the model using them.
% In the second step, we update the outcome and ITE propagation terms using a mini-batch consisting of $b_2$ instance pairs. 
In the first step, we sample a mini-batch consisting of $b_1$ labeled instances to approximate the supervised loss (\ref{eq:supervised}).
In the second step, we compute the outcome propagation term and the ITE propagation terms using a mini-batch consisting of $b_2$ instance pairs. 
Note that in order to make the model more flexible, we can employ different regularization parameters for the treatment outcomes and the control outcomes. 
The $b_1$ and $b_2$ are considered as hyper-parameters; the details are described in Section~\ref{sec:experiments}.
In practice, we optimize only the supervised loss for the first several epochs, and decrease the strength of regularization as training proceeds, in order to guide efficient training~\cite{weston2012deep}.  
\begin{algorithm}[tb]
  \caption{Counterfactual propagation}\label{alg:train}
  \SetAlgoLined
  % \State {\bf Input:} \mathcal{P},\mathcal{I},\mathcal{R}, t_{ui}>0\\
  % \State {\bf Output:} the latent factors, {\bf p}_u, {\bf q}_i\\
%   {\bf Initialize:} $\Theta$\\
%   {\# Estimate the dwell-time distribution for each user}\\
 {\bf Input:} labeled instances~$\{(\x_i, t_i, y^{t_i}_i)\}^{N}_{i=1}$, unlabeled instances~$\{(\x_i)\}^{N+M}_{i=N+1}$, a similarity matrix $w=(w_{ij})$, and mini-batch sizes $b_1, b_2$.\\
 {\bf Output:} estimated outcome(s) for each treatment $\hat{y}^{1}_i$ and/or $\hat{y}^{0}_i$ using Eq.(~\ref{eq:tarnet}~).   \\
%   \For{instance $i$}{
%   neighbors_{i}$\leftarrow$Find nearest neighbors top-k\\
%     }
%   {\# BPR framework}
%   {\# Train a BPR model}\\
  \While{not converged}{
%   {\# Minimizing supervised loss }\\
  {\# Approximating the supervised loss }\\
    Sample $b_1$ instances $\{(\x_i, t_i, y_i)\}$ from the labeled instances\\
%     Update $f$ to minimize the mini-batch approximation of the supervised loss (\ref{eq:supervised})\\
%     {\# Minimizing propagation terms}\\
%         Sample $b_2$ pairs of instances~$\{(\x_i, \x_j)\}$\\
%     Update $f$ to minimize the mini-batch approximation of the propagation terms $\lambda_\text{o}L_\text{o}$\\
%     Sample $b_2$ pairs of instances~$\{(\x_i, \x_j)\}$\\
%     Update $f$ to minimize the mini-batch approximation of the propagation terms $\lambda_\text{e}L_\text{e}$\\
Compute the supervised loss (\ref{eq:supervised}) for the $b_1$ instances\\
 {\# Approximating propagation terms}\\
  Sample $b_2$ pairs of instances~$\{(\x_i, \x_j)\}$\\
 Compute the outcome propagation terms $\lambda_\text{o}L_\text{o}$ for $b_2$ pairs of instances\\
   Sample $b_2$ pairs of instances~$\{(\x_i, \x_j)\}$\\
 Compute the ITE propagation terms $\lambda_\text{e}L_\text{e}$ of $b_2$ pairs of instances\\
 Update the parameters to minimize $L_s+\lambda_{o}L_o+\lambda_{e}L_e$ for the sampled instances\\
 
    }
\end{algorithm}


\section{Experiments\label{sec:experiments}}
%We discuss experiments in order to validate our method. 
% \harada{
We test the effectiveness of the proposed semi-supervised ITE estimation method in comparison with various supervised baseline methods, especially when the available labeled data are strictly limited. 
We first conduct experiments using two semi-synthetic datasets based on public real datasets. 
We also design some experiments varying the magnitude of noise on outcomes to explore how the noisy outcomes affect the proposed method.
% Our implementation is available on Github\footnote{\url{https://anonymous.4open.science/r/6ad51d8e-c33c-442f-9a31-f603612bde44/} (Note that this is an anonymous account.)}
%\harada{Our implementation will be available if the paper is accepted.}

% We also conduct experiments using two synthetic datasets. 
% We also conduct experiments using two synthetic datasets. 
%main interests lies in the following three research questions: 

% \vspace{2mm}

% {\bf RQ1:} Does the semi-supervised learning improve the ITE estimation accuracy when the number of labeled instance is small? 

% \vspace{2mm}
% {\bf RQ2:} Is the proposed model robust to the choice of the regularization hyperparameters? 

% \noindent
% {\bf RQ2:} Is the proposed model robust to noisy outcomes?
%\footnote{\harada{Our Implementation of this study is available ..}}
% In both setting, our goal is to explore how the performances change in response to the size of labeled data and validate the advantage of the proposed method.}

% \kashima{Sometimes it is nice to first give research questions we want to answer through experiments. For example, we can say like ``We conducted experiments using .... to investigate the following questions: (1) does the semi-supervised learning improves the estimation accuracy when the number of labeled instance is small, (2) .... "}

\subsection{Datasets}

Owing to the counterfactual nature of ITE estimation, we rarely access real-world datasets including ground truth ITEs, and therefore cannot directly evaluate ITE estimation methods like the standard supervised learning methods using cross-validation.
Therefore, following the existing work~\cite{johansson2016learning}, we employ two semi-synthetics datasets whose counterfactual outcomes are generated through simulations. 
Refer to the original papers for the details on outcome generations~\cite{hill2011bayesian,johansson2016learning}.


\subsubsection{News dataset} is a dataset including opinions of media consumers for news articles~\cite{johansson2016learning}. 
It contains $5{,}000$ news articles and outcomes generated from the NY Times corpus\footnote{\url{https://archive.ics.uci.edu/ml/datasets/Bag+of+Words}}.
Each article is consumed on desktop~($t=0$)~or mobile~($t=1$) and it is assumed that media consumers prefer to read some articles on mobile than desktop. 
% Each article is represented using bag of words and the size of the vocabulary is $3{,}477$. 
Each article is generated by a topic model and represented in the bag-of-words representation. The size of the vocabulary is $3{,}477$. 
% \harada{We use the subset of dataset employed in the previous work}~\cite{johansson2016learning}.


\subsubsection{IHDP dataset} is a dataset created by randomized experiments called the Infant Health and Development Program (IHDP)~\cite{hill2011bayesian} to examine the effect of special child care on future test scores. 
It contains the results of $747$ subjects~($139$ treated subjects and $608$ control subjects)~with $25$ covariates related to infants and their mothers.
Following the existing studies~\cite{johansson2016learning,shalit2017estimating}, the ground-truth counterfactual outcomes are simulated using the NPCI package~\cite{dorie2016npci}. 


\subsubsection{Synthetic dataset} is a synthetically generated dataset for this study. 
We generate $1{,}000$ instances that have eight covariates sampled from $\mathcal{N}({\bf 0}^{8\times1},0.5\times(\Sigma+\Sigma^{\top} ))$,
where $\Sigma $ is sampled from $\sim\mathcal{U}((-1,1)^{8\times8})$.
The treatment $t$ is sampled from $\text{Bern}(\sigma({\bf w}_t^{\top}{\bf x}+\epsilon_t))$, where $\sigma(\cdot)=\frac{1}{1+\exp(-\cdot)}, {\bf w}_t\sim\mathcal{U}((-1,1)^{8\times1}), \epsilon_t\sim\mathcal{N}(0, 0.1)$. 
The treatment outcome and control outcome are generated as  $y^{1}|\x\sim \sin({\bf w}^{\top}_y \x) +c\epsilon_{y}$ and $y^{0}|\x\sim \cos({\bf w}^{\top}_y \x) +c\epsilon_{y}$,
respectively, where ${\bf w} \sim\mathcal{U}((-1,1)^{8\times1})$ and $ \epsilon_y\sim\mathcal{N}({\bf 0}, 1)$. 


% \harada{\kashima{Any reason to use a subset? $\rightarrow$}
% We use the subset of dataset employed in the previous work}~\cite{shalit2017estimating}.\haradap{No. I will delete this sentence.}

% \subsubsection{Synthetic dataset}\harada{write systems}. 




% \subsection{RQ1: Does the semi-supervised learning improve the ITE estimation accuracy when the number of labeled instance is small? }

\subsection{Experimental settings}
% \subsubsection{Experimental settings}
Since we are particularly interested in the situation when the available labeled data are strictly limited, we split the data into a training dataset, validation dataset, and a test dataset by limiting the size of the training data. 
We change the ratio of the training to investigate the performance; we use $10\%, 5\%,$ and $1\%$ of the whole data from the News dataset, and use $40\%, 20\%$, and $10\%$ of those from the IHDP dataset for the training datasets. 
% On IHDP dataset, we hold $50\%$ and $10\%$ of the whole data for test and validation, respectively. 
The rest $80\%$ and $10\%$ of the whole News data are used for test and validation, respectively.
Similarly, $50\%$ and $10\%$ of the whole IHDP dataset are used for test and validation, respectively.
% On the News dataset, we hold $80\%$ and $10\%$ for test and validation. 
% We report the average results of $10$ trials on the News dataset and $50$ trials on the IHDP dataset. 
We report the average results of $10$ trials on the News dataset, $50$ trials on the IHDP dataset, and $10$ trials on the Synthetic datasets.


In addition to the evaluation under labeled data scarcity, we also test the robustness against label noises. 
As pointed out in previous studies, noisy labels in training data can severely deteriorate predictive performance, especially in semi-supervised learning.  
Following the previous  work~\cite{hill2011bayesian,johansson2016learning}, we add the noise $\epsilon\sim\mathcal{N}(0, c^2)$ to the observed outcomes in the training data, where $c \in \{1, 3, 5, 7, 9\}$.
%Note that in both datasets, the outcomes employed in the previous studies~\cite{hill2011bayesian,johansson2016learning,shalit2017estimating} and this study are generated as $c=1$ unless mentioned explicitly.  
In this evaluation, we use $1\%$ of the whole data as the training data for the News dataset and $10\%$  for the IHDP dataset, respectively, since we are mainly interested in label-scarce situations.



The hyper-parameters are tuned based on the prediction loss using the  observed outcomes on the validation data. 
We calculate the similarities between the instances by using the Gaussian kernel; 
we select $\sigma^{2}$ from $\{5\times10^{-3}, 1\times10^{-3},\ldots , 1\times10^{2}, 5\times10^{2}\}$, 
and select $\lambda_o$ and $\lambda_e$ from $\{1\times10^{-3}, 1\times10^{-2}, \ldots, 1\times10^{2}\}$. 
Because the scales of treatment outcomes and control outcomes are not always the same, we found scaling the regularization terms according to them is beneficial; 
% incorporating them in setting of the regularization hyperparameters is beneficial according to our observations;
specifically, we scale the regularization terms with respect to the treatment outcomes, the control outcomes, and the treatment effects by $\alpha={1}/{\sigma_{y^1}^{2}}, \beta={1}/{\sigma_{y^0}^{2}}$, and $\gamma={1}/{(\sigma_{y^1}^{2}+\sigma_{y^0}^{2})}$, respectively.
%, and use them for the regularization hyperparameters \harada{as $\alpha\lambda_o$, $\beta\lambda_o$, $\gamma\lambda_e$ for treatment outcome, control outcome, and treatment effect regularization, respectively}.
We apply principal component analysis to reduce the input dimensions before applying the Gaussian kernel; we select the number of dimensions from $\{2, 4, 6, 8, 16, 32, 64\}$. 
The learning rate is set to $1\times10^{-3}$ and the mini-batch sizes $b_1, b_2$ are chosen from $\{4,8,16,32\}$. 

As the evaluation metrics, we report the {\it Precision in Estimation of Heterogeneous Effect~(PEHE)~} used in the previous research~\cite{hill2011bayesian}.
PEHE is the estimation error of individual  treatment effects, and is defined as 
\[
\epsilon_{\rm PEHE}=\frac{1}{N+M}\sum_{i=1}^{N+M}(\tau_i - \hat{\tau}_i)^2.
\]
Following the previous studies~\cite{shalit2017estimating,yao2018representation}, we evaluate the predictive performance for labeled instances and unlabeled instances separately. 
% The first task is to estimate ITE for instances which are included in the train dataset~(which we refer to as `within-sample'). 
% Although we observe factual outcomes of instances in the training data,  this setting still evaluate predictive performance because we cannot access to the counterfactual outcomes. 
% The other task is to estimate ITE for instances which are not included in the training data~(which we call `without-sample'). 
% The first task is conducted using the training and validation dataset in the same way as existing work~\cite{shalit2017estimating}. 
%The first task is to estimate the ITEs for the labeled instances. 
Note that, although we observe the factual outcomes of the labeled data, their true ITEs are still unknown because we cannot observe their counterfactual outcomes.
%Note that we do not know them because only the factual outcomes are given in the labeled data.
%of instances in the training data,  this setting still evaluate predictive performance because we cannot access to the counterfactual outcomes. 
%The other task is to estimate the ITEs for unlabeled instances.
%The first task is conducted using the training and validation datasets following the procedure proposed in the previous work~\cite{shalit2017estimating}. 

\begin{table*}[tbp]
\caption{The performance comparison of different methods on News dataset. The $^\dagger$ indicates that our proposed method (CP) performs statistically significantly better than the baselines by the paired $t$-test ($p<0.05)$. The bold results indicate the best results in terms of the average.}\label{table:news}
% \scalebox{0.919}[0.92]
{
%   \begin{center}
\centering

%     \begin{tabular}{c}
      % 1
%       \begin{minipage}{\hsize}
\begin{tabular} {ccccccccc }
% \caption{Experimental results}

\toprule[2pt]
          $ \sqrt{\epsilon_{\rm PEHE}} $         & \multicolumn{2}{c}{ \bf{News} 1\%}        &              & \multicolumn{2}{c}{\bf{News} 5\%}& &  \multicolumn{2}{c}{\bf{News} 10\%}      \\[3pt]\cmidrule{2-3}\cmidrule{5-6}\cmidrule{8-9}
Method                      & labeled  & unlabeled & \multicolumn{1}{l}{} & labeled & unlabeled & \multicolumn{1}{l}{} & labeled & unlabeled\\[3pt] \cmidrule{1-9}
Ridge-1   &  $^\dagger4.494_{\pm{1.116}}$  &   $^\dagger4.304_{\pm{0.988}}$         & \multicolumn{1}{l}{} &  $^\dagger4.666_{\pm{1.0578}}$ &  $^\dagger3.951_{\pm{0.954}}$ &   & $^\dagger4.464_{\pm{1.082}}$   & $^\dagger3.607_{\pm{0.943}}$      \\[3pt]
Ridge-2   &  $2.914_{\pm{0.797}}$ &  $^\dagger2.969_{\pm{0.814}}$  & \multicolumn{1}{l}{} & $^\dagger2.519_{\pm{0.586}}$ & $^\dagger2.664_{\pm{0.614}}$ &  & $^\dagger2.560_{\pm{0.558}}$ & $^\dagger2.862_{\pm{0.621}}$      \\[3pt]

Lasso-1   &        $^\dagger4.464_{\pm{1.082}}$   &   $^\dagger3.607_{\pm{0.943}}$   & \multicolumn{1}{l}{} &    $^\dagger4.466_{\pm{1.058}}$    &     $^\dagger3.367_{\pm{0.985}}$    &  &   $^\dagger4.464_{\pm{1.0822}}$& $^\dagger3.330_{\pm{0.984}}$      \\[3pt]
Lasso-2   &        $^\dagger3.344_{\pm{1.022}}$   &   $^\dagger3.476_{\pm{1.038}}$   & \multicolumn{1}{l}{} &    $^\dagger2.568_{\pm{0.714}}$   &  $^\dagger2.848_{\pm{0.751}}$ &  &  $^\dagger2.269_{\pm{0.628}}$  &  $^\dagger2.616_{\pm{0.663}}$      \\[3pt]

\cmidrule{1-9}
\multicolumn{1}{c}{$k$NN} &    $^\dagger3.678_{\pm{1.250}}$   &   $^\dagger3.677_{\pm{1.246}}$  & \multicolumn{1}{l}{} &   $^\dagger3.351_{\pm{1.004}}$&   $^\dagger3.434_{\pm{1.018}}$  &  &  $^\dagger3.130_{\pm{0.752}}$ & $^\dagger3.294_{\pm{0.766}}$      \\[3pt]
\multicolumn{1}{c}{PSM}    &  $^\dagger3.713_{\pm{1.149}}$  &  $^\dagger3.662_{\pm{1.127}}$ & \multicolumn{1}{l}{} &   $^\dagger3.363_{\pm{0.901}}$  &  $^\dagger3.500_{\pm{0.961}}$  &  &   $^\dagger3.260_{\pm{0.734}}$& $^\dagger3.526_{\pm{0.832}}$      \\[3pt]
\cmidrule{1-9}
\multicolumn{1}{c}{RF}    &  $^\dagger4.494_{\pm{1.116}}$  &    $^\dagger3.691_{\pm{0.878}}$   & \multicolumn{1}{l}{} &    $^\dagger4.466_{\pm{1.058}}$   &  $^\dagger2.975_{\pm{0.874}}$    &  &  $^\dagger4.464_{\pm{1.082}}$ & $^\dagger2.657_{\pm{0.682}}$     \\[3pt]
\multicolumn{1}{c}{CF}    &  $^\dagger3.691_{\pm{1.082}}$  &    $^\dagger3.607_{\pm{0.943}}$   & \multicolumn{1}{l}{} &  $^\dagger3.196_{\pm{0.901}}$ & $^\dagger3.215_{\pm{0.910}}$ &  & $^\dagger3.101_{\pm{0.806}}$ & $^\dagger3.129_{\pm{0.818}}$ \\[3pt]
\cmidrule{1-9}
\multicolumn{1}{c}{TARNET}    &        $^\dagger3.166_{\pm{0.742}}$       &       $^\ddagger3.160_{\pm{0.722}}$         & \multicolumn{1}{l}{} &   $^\dagger2.670_{\pm{0.796}}$  &   $^\dagger2.666_{\pm{0.773}}$      &  &   $^\dagger2.589_{\pm{0.894}}$ & $^\dagger2.598_{\pm{0.869}}$      \\[3pt]
\multicolumn{1}{c}{CFR}    &        $2.908_{\pm{0.752}}$       &       $2.925_{\pm{0.746}}$         & \multicolumn{1}{l}{} &   $^\dagger2.590_{\pm{0.772}}$   &  $^\dagger2.546_{\pm{0.796}}$   &  &   $^\dagger2.570_{\pm{0.519}}$& $^\dagger2.451_{\pm{0.547}}$      \\[3pt]
\cmidrule{1-9}
CP (proposed) &        $\bf 2.844_{\pm{0.683}}$       &       $\bf 2.823_{\pm{0.656}}$        & \multicolumn{1}{l}{} &     $\bf 2.310_{\pm{0.430}}$          & $\bf 2.446_{\pm{0.471}}$       &  &   $\bf 2.003_{\pm{0.393}} $& $\bf 2.153_{\pm{0.436}}$      \\[3pt]
% \cmidrule{1-9}
% Improvement\\
\bottomrule[2pt]
    \end{tabular}
    }
\end{table*}
\begin{table*}[tbp]
\caption{The performance comparison of different methods on IHDP dataset. The $^\dagger$ indicates that our proposed method (CP) performs statistically significantly better than the baselines by the paired $t$-test ($p<0.05$). The bold results indicate the best results in terms of average.}\label{table:ihdp}
\scalebox{0.975}[0.975]
{
\centering
\begin{tabular} {ccccccccc }
\toprule[2pt]
       $ \sqrt{\epsilon_{\rm PEHE}} $         & \multicolumn{2}{c}{ \bf{IHDP} 10\%}        &              & \multicolumn{2}{c}{\bf{IHDP} 20\%}& &  \multicolumn{2}{c}{\bf{IHDP} 40\%}      \\[3pt]\cmidrule{2-3}\cmidrule{5-6}\cmidrule{8-9}
Method                      & labeled & unlabeled & \multicolumn{1}{l}{} & labeled & unlabeled & \multicolumn{1}{l}{} & labeled & unlabeled\\[3pt] \cmidrule{1-9}
Ridge-1   &        $^\dagger5.484_{\pm{8.825}}$ &  $^\dagger5.696_{\pm{7.328}}$ & \multicolumn{1}{l}{} &   $^\dagger5.067_{\pm{8.337}}$  &  $^\dagger4.692_{\pm{6.943}}$       &   &$^\dagger4.80_{\pm{8.022}}$   & $^\dagger4.448_{\pm{6.874}}$      \\[3pt]
Ridge-2   &        $^\dagger3.426_{\pm{5.692}}$      &       $^\dagger3.357_{\pm{5.177}}$       & \multicolumn{1}{l}{} &  $^\dagger2.918_{\pm{4.874}}$  &    $^\dagger2.918_{\pm{4.730}}$  &   &  $^\dagger2.605_{\pm{4.314}}$ & $^\dagger2.639_{\pm{4.496}}$      \\[3pt]

Lasso-1   &        $^\dagger6.685_{\pm{10.655}}$       &       $^\dagger6.408_{\pm{9.900}}$         & \multicolumn{1}{l}{} &        $^\dagger6.435_{\pm{10.147}}$       &  $^\dagger6.2446_{\pm{9.639}}$       &     &   $^\dagger6.338_{\pm{9.704}}$ &  $^\dagger6.223_{\pm{9.596}}$  \\[3pt]
Lasso-2   &        $^\dagger3.118_{\pm{5.204}}$       &       $^\ddagger3.292_{\pm{5.725}}$         & \multicolumn{1}{l}{} &         $^\dagger2.684_{\pm{4.428}}$    &  $^\dagger2.789_{\pm{4.731}}$       &  &   $^\dagger2.512_{\pm{4.075}}$& $^\dagger2.571_{\pm{4.379}}$      \\[3pt]
\cmidrule{1-9}
\multicolumn{1}{c}{$k$NN} &        $^\dagger4.457_{\pm{6.957}}$  &  $^\dagger4.603_{\pm{6.629}}$         & \multicolumn{1}{l}{} &      $^\dagger4.023_{\pm{6.193}}$     &    $^\dagger4.370_{\pm{6.244}}$     &  &   $^\dagger3.623_{\pm{5.316}}$ & $^\dagger4.109_{\pm{5.936}}$      \\[3pt]
\multicolumn{1}{c}{PSM}    &        $^\dagger6.506_{\pm{10.077}}$       &       $^\dagger6.982_{\pm{10.672}}$         & \multicolumn{1}{l}{} &   $^\dagger6.277_{\pm{9.708}}$     &    $^\dagger7.209_{\pm{11.077}}$    &  &   $^\dagger6.065_{\pm{9.362}}$ & $^\dagger7.181_{\pm{9.362}}$      \\[3pt]
\cmidrule{1-9}
\multicolumn{1}{c}{RF}    &        $^\dagger6.924_{\pm{10.620}}$       &       $^\dagger5.356_{\pm{8.790}}$         & \multicolumn{1}{l}{} &   $^\dagger6.854_{\pm{10.471}}$      &    $^\dagger4.845_{\pm{8.241}}$     &  &   $ ^\dagger6.928_{\pm{10.396}}$ & $^\dagger4.549_{\pm{7.822}}$\\[3pt]
\multicolumn{1}{c}{CF}    &        $^\dagger5.389_{\pm{8.736}}$       &       $^\dagger5.255_{\pm{8.070}}$         & \multicolumn{1}{l}{} &   $^\dagger4.939_{\pm{7.762}}$      &    $^\dagger4.955_{\pm{7.503}}$   &  &   $^\dagger4.611_{\pm{7.149}}$  &  $^\dagger4.764_{\pm{7.448}}$\\[3pt]
\cmidrule{1-9}
\multicolumn{1}{c}{TARNET}    &        $^\dagger3.827_{\pm{5.315}}$       &       $^\dagger3.664_{\pm{4.888}}$         & \multicolumn{1}{l}{} &   $^\dagger2.770_{\pm{3.617}}$  &    $^\dagger2.770_{\pm{3.542}}$   &     &   $^\dagger2.005_{\pm{2.447}}$& $^\dagger2.267_{\pm{2.825}}$      \\[3pt]
\multicolumn{1}{c}{CFR}    &        $^\dagger3.461_{\pm{5.1444}}$       &       $^\dagger3.292_{\pm{4.619}}$         & \multicolumn{1}{l}{} &  $^\dagger2.381_{\pm{3.126}}$ &    $^\dagger2.403_{\pm{3.080}}$      &  &   $1.572_{\pm{1.937}}$ & $1.815_{\pm{2.204}}	$   	 \\[3pt]
\cmidrule{1-9}
CP~(proposed) &       $\bf 2.427_{\pm{3.189}}$ & $\bf 2.652_{\pm{3.469}}$        & \multicolumn{1}{l}{} &     $\bf 1.686_{\pm{1.838}}$          & $\bf 1.961_{\pm{2.343}}$       &  &   $\bf 1.299_{\pm{1.001}} $& $\bf 1.485_{\pm{1.433}}$      \\[3pt]
\bottomrule[2pt]
    \end{tabular}
    }
\end{table*}

\begin{table*}[tbp]
\caption{The performance comparison of different methods on Synthetic dataset. The $^\dagger$ indicates that our proposed method (CP) performs statistically significantly better than the baselines by the paired $t$-test ($p<0.05$). The bold results indicate the best results in terms of average.}\label{table:synthetic}
\scalebox{0.975}[0.975]
{
\centering
\begin{tabular} {ccccccccc }
\toprule[2pt]
       $ \sqrt{\epsilon_{\rm PEHE}} $         & \multicolumn{2}{c}{ \bf{Synthetic} 10\%}        &              & \multicolumn{2}{c}{\bf{Synthetic} 20\%}& &  \multicolumn{2}{c}{\bf{Synthetic} 40\%}      \\[3pt]\cmidrule{2-3}\cmidrule{5-6}\cmidrule{8-9}
Method                      & labeled & unlabeled & \multicolumn{1}{l}{} & labeled & unlabeled & \multicolumn{1}{l}{} & labeled & unlabeled\\[3pt] \cmidrule{1-9}
Ridge-1   &        $^\dagger0.966_{\pm{0.078}}$ &  $^\dagger0.960_{\pm{0.78}}$ & \multicolumn{1}{l}{} &   $^\dagger0.965_{\pm{0.072}}$  &  $^\dagger0.960_{\pm{0.79}}$       &   & $^\dagger0.959_{\pm{0.062}}$ &  $^\dagger0.949_{\pm{0.77}}$      \\[3pt]
Ridge-2   &        $^\dagger0.890_{\pm{0.180}}$  &  $^\dagger0.950_{\pm{0.193}}$       & \multicolumn{1}{l}{}
&  $^\dagger0.858_{\pm{0.164}}$  &    $^\dagger0.881_{\pm{0.193}}$  &   &  $^\dagger0.836_{\pm{0.137}}$ & $^\dagger0.855_{\pm{0.175}}$      \\[3pt]

Lasso-1   &        $^\dagger0.965_{\pm{0.077}}$       &       $^\dagger0.959_{\pm{0.075}}$         & \multicolumn{1}{l}{} &        $^\dagger0.964_{\pm{0.069}}$       &  $^\dagger0.951_{\pm{0.074}}$       &     &   $^\dagger0.959_{\pm{0.061}}$ &  $^\dagger0.949_{\pm{0.074}}$  \\[3pt]

Lasso-2   &        $^\dagger0.862_{\pm{0.121}}$       &       $^\ddagger0.887_{\pm{0.142}}$         & \multicolumn{1}{l}{} &         $^\dagger0.877_{\pm{0.111}}$    &  $^\dagger0.876_{\pm{0.126}}$       &  &   $^\dagger0.872_{\pm{0.104}}$& $^\dagger0.871_{\pm{0.124}}$     \\[3pt]\cmidrule{1-9}

\multicolumn{1}{c}{$k$NN} &        $^\dagger0.787_{\pm{0.104}}$  &  $^\dagger0.842_{\pm{0.126}}$         & \multicolumn{1}{l}{} &      $^\dagger0.703_{\pm{0.107}}$     &    $^\dagger0.766_{\pm{0.134}}$     &  &   $^\dagger0.643_{\pm{0.096}}$ & $^\dagger0.695_{\pm{5.936}}$  \\[3pt]

\multicolumn{1}{c}{PSM}    &        $^\dagger0.972_{\pm{0.085}}$       &       $^\dagger0.989_{\pm{0.093}}$         & \multicolumn{1}{l}{} &   $^\dagger0.957_{\pm{0.769}}$     &    $^\dagger0.980_{\pm{0.088}}$    &  &   $^\dagger0.940_{\pm{0.060}}$ & $^\dagger0.975_{\pm{0.080}}$      \\[3pt]
\cmidrule{1-9}

\multicolumn{1}{c}{RF}    &        $^\dagger0.987_{\pm{0.0543}}$       &       $^\dagger0.902_{\pm{0.119}}$         & \multicolumn{1}{l}{} &   $^\dagger1.001_{\pm{0.033}}$      &    $^\dagger0.839_{\pm{0.144}}$     &  &   $ ^\dagger1.014_{\pm{0.041}}$ & $^\dagger0.804_{\pm{0.158}}$\\[3pt]
\multicolumn{1}{c}{CF}    &        $^\dagger0.852_{\pm{0.082}}$       &       $^\dagger0.890_{\pm{0.099}}$         & \multicolumn{1}{l}{} &   $^\dagger0.819_{\pm{0.082}}$      &    $^\dagger0.855_{\pm{0.106}}$   &  &   $^\dagger0.786_{\pm{0.076}}$  &  $^\dagger0.826_{\pm{0.102}}$\\[3pt]
\cmidrule{1-9}


\multicolumn{1}{c}{TARNET}    &        $^\dagger0.522_{\pm{0.165}}$       &       $^\dagger0.556_{\pm{0.193}}$         & \multicolumn{1}{l}{} &   $\dagger0.326_{\pm{0.821}}$  &    $\dagger0.363_{\pm{0.101}}$   &     &   $^\dagger0.225_{\pm{0.026}}$& $^\dagger0.260_{\pm{0.041}}$      \\[3pt]
\multicolumn{1}{c}{CFR}    &        $^\dagger0.521_{\pm{0.166}}$       &       $^\dagger0.561_{\pm{0.193}}$         & \multicolumn{1}{l}{} &  $\dagger0.311_{\pm{0.072}}$ &    $\dagger0.348_{\pm{0.861}}$      &  &   $\bf{0.215_{\pm{0.019}}}$ & $\dagger0.258_{\pm{0.036}}	$   	 \\[3pt]
\cmidrule{1-9}

CP~(proposed) &       $\bf 0.307_{\pm{0.125}}$ & $\bf 0.331_{\pm{0.149}}$        & \multicolumn{1}{l}{} &     $\bf 0.249_{\pm{0.063}}$          & $\bf 0.259_{\pm{0.086}}$       &  &   $0.229_{\pm{0.031}} $& ${\bf 0.211_{\pm{0.049}}}$      \\[3pt]

\bottomrule[2pt]
    \end{tabular}
    }
\end{table*}





\subsection{Baselines}
% \subsubsection{Baselines}
We compare the proposed method with several existing supervised ITE estimation approaches.
(i) Linear regression~(Ridge, Lasso) is the ordinary linear regression models with ridge regularization or lasso regularization. 
We consider two variants: one that includes the treatment as a feature~(denoted by `Ridge-$1$' and `Lasso-$1$'),
and the other with two separated models for treatment and control~(denoted by `Ridge-$2$' and `Lasso-$2$').  
% \vspace{1mm}
% \noindent
% {\bf LinearRegression-2:} is separated ordinary linear regression model for treatment and control.
% \vspace{1mm}
% \noindent
(ii) $k$-nearest neighbors~($k$NN) is a matching-based method that predicts the outcomes using nearby instances.
% \vspace{1mm}
% \noindent
(iii) Propensity score matching with logistic regression~(PSM)~\cite{rosenbaum1983central} is a matching-based method using the propensity score estimated by a logistic regression model.
% \vspace{1mm}
% \noindent
% {\bf Random forest~(RF)}~\cite{breiman2001random} and its causal extension called 
% {\bf Causal forest~(CF)}~\cite{wager2018estimation} are also tested.
We also compared the proposed method with tree models such as  (iv) random forest~(RF)~\cite{breiman2001random} and its causal extension called (v) causal forest~(CF)~\cite{wager2018estimation}. 
In CF, trees are trained to predict propensity score and leaves are used to predict treatment effects. 
% \kashima{Include TARNet. I am not sure the difference between TARNet and CFR. Both seem intoduced in Shalit2017.???} \haradap{Actually both are included in shalit2017. TARNet + balancing = CFR. }
% (vi) TARNet~\cite{shalit2017estimating} and (vii) Counterfactual regression~(CFR-Wass)~\cite{shalit2017estimating} are state-of-the-art DNN models based on balanced representations between treatment and control instances. We use the Wasserstein model. 
% \kashima{You can just call CFR (instead of CFRWass)?}
(vi) TARNet~\cite{shalit2017estimating}  is a deep neural network model that has shared layers for representation learning and different layers for outcome prediction for treatment and control instances. 
%See Section~\ref{sec:proposed} for more details.  
(vii) Counterfactual regression~(CFR)~\cite{shalit2017estimating} is a  state-of-the-art deep neural network model based on balanced representations between treatment and control instances. We use the Wasserstein distance. 



\subsection{Results and discussions}
% \kashima{RQs should be the same as the ones we gave in the beginning of this section.}
% \subsubsection{Results and discussions}
% \kashima{Does this section test performance FOR SMALL LABELED DATA?}\haradap{Yes}
%We discuss the experimental results.
% \harada{In this subsection, we mainly discuss how the predictive performance changes according to the size of labeled data and the magnitude of noise. First we explain experimental results varying the size of labeled data and sensitivity to hyper-parameters which control the strength of propagation. Then we analyze the experimental results on noisy outcomes.}
We discuss the performance of the proposed method compared with the baselines by changing the size of labeled datasets, and then investigate the robustness against the label noises. 

We first see the experimental results for different sizes of labeled datasets and sensitivity to the choice of the hyper-parameters that control the strength of label propagation. 
% 全体的な傾向、提案手法良い。NEWSで特に良い
% Tables \ref{table:news} and \ref{table:ihdp} show the PEHE values by different methods for the News dataset and the IHDP dataset, respectively. 
% Overall, our proposed method exhibits the best ITE estimation performance for both labeled and unlabeled data in both of the datasets; the advantage is more significant in the News dataset. 
%\harada{
Tables \ref{table:news}, \ref{table:ihdp}, and \ref{table:synthetic} show the PEHE values by different methods for the News dataset, the IHDP dataset, and the Synthetic dataset, respectively. 
Overall, our proposed method exhibits the best ITE estimation performance for both labeled and unlabeled data in all of the three datasets.
%Among the two semi-synthetic datasets (News and IHDP), the advantage is more significant in the News dataset. 
%}
In general, the performance gain by the proposed method is larger on labeled data than on unlabeled data. 

% NEWSでの話
%On the News dataset, the proposed method achieves the best performance in terms of PEHE for both the labeled and unlabeled instances  settings. 
The News dataset is a relatively high-dimensional dataset represented using a bag of words.
%seems to be relatively easier for defining  appropriate similarity measures. 
The two-model methods such as Ridge-2 and Lasso-2 perform well in spite of their simplicity, and in terms of regularization types, the Lasso-based methods perform relatively better due to the high-dimensional nature of the dataset. 

% IHDPの話 と やや精度が落ちる理由
The proposed method also performs the best in the IHDP dataset; however, the performance gain is rather moderate, as shown by the no statistical significance against CFR~\cite{shalit2017estimating} with the largest $40\%$-labeled data, which is the most powerful baseline method.
The reason for the moderate improvements is probably because of the difficulty in defining appropriate similarities among instances, because the IHDP dataset has various types of features including continuous variables and discrete variables.
%The traditional baselines including Ridge1, Lasso1, kNN matching, and the tree based models do not degrade their predictive accuracy significantly due to the scarcity of labeled data, however their predictive capability still limit their performance even when we increase the labeled data. 
% \kashima{Traditional baselines such as Ridge1, Lasso1, kNN matching, and the tree based models show limited performance; }
The traditional baselines such as Ridge-1, Lasso-1, $k$-NN matching, and the tree-based models show limited performance;
in contrast, the deep learning based methods such as TARNet and CFR demonstrate remarkable performance. 
%However, these methods seem to be vulnerable to the decrease of labeled data.  

% Syntheticの話
%\harada{
The proposed method again achieves the best performance in the Synthetic dataset. 
Since the outcomes are generated by a non-linear function, the linear regression methods are not capable of capturing the non-linear behaviour, and therefore show the limited performance. 
Though TARNet and CFR show the excellent results particularly when using $40$\% of the whole dataset, they suffer from the scarcity of labeled data, and significantly degrade their performance. 
On the other hand, the proposed method is rather robust to the lack of labeled data.
%}



\begin{table}[tbp]
\caption{Investigation of the contributions by the outcome propagation and the ITE propagation in the proposed method.
% The upper table shows the results for the News dataset, and the lower for the IHDP dataset.
The top table shows the results for the News dataset, the middle one for the IHDP, and the bottom one for the Synthetic dataset.
%on the News dataset~(upper)~and the IHDP dataset~(lower). 
The $\lambda_o=0$ and $\lambda_e=0$ indicate the proposed method (CP) with only the ITE propagation and the outcome propagation, respectively,
The $^\dagger$ indicates that our proposed method (CP) performs statistically significantly better than the baselines by the paired $t$-test ($p<0.05)$. 
The bold numbers indicate the best results in terms of the average.}\label{table:proposed_comparison} 
% On IHDP dataset, the outcome regularizations seem to play a more important role than the treatment effect resgularization. On News dataset, when smaller size of labeled data, treatment effect regularization shows slightly better effects in terms of average.}}
\vspace{2mm}
  \centering
\begin{tabular} {ccccccccc }
\toprule[2pt]
        $ \sqrt{\epsilon_{\rm PEHE}} $         & \multicolumn{2}{c}{ \bf{News} 1\%}        &              & \multicolumn{2}{c}{\bf{News} 5\%}& &  \multicolumn{2}{c}{\bf{News} 10\%}      \\[3pt]\cmidrule{2-3}\cmidrule{5-6}\cmidrule{8-9}
Method                      & labeled  & unlabeled & \multicolumn{1}{l}{} & labeled & unlabeled & \multicolumn{1}{l}{} & labeled & unlabeled\\[3pt] \cmidrule{1-9}

CP~($\lambda_o=0$)  &  $\bf 2.812_{\pm{651}}$ & $\bf 2.806_{\pm{0.598}}$  & \multicolumn{1}{l}{} &  $^\dagger2.527_{\pm{0.474}}$  & $\dagger2.531_{\pm{0.523}}$  &     &  $^\dagger2.400_{\pm{0.347}}$ & $^\dagger2.410_{\pm{0.450}}$  \\[3pt]

CP~($\lambda_e=0$) &        $2.879_{\pm{0667}}$       &       $2.885_{\pm{0.609}}$         & \multicolumn{1}{l}{} &     $ 2.351_{\pm{0.450}}$   &   $2.483_{\pm{0.481}}$      &       &  $\bf 1.996_{\pm{0.338}}$ & $2.221_{\pm{0.455}}$   \\[3pt]

CP &        $2.844_{\pm{0.683}}$       &       $2.823_{\pm{0.656}}$        & \multicolumn{1}{l}{} &     $\bf 2.310_{\pm{0.430}}$          & $\bf 2.446_{\pm{0.471}}$       &  &   $2.003_{\pm{0.393}} $& $\bf 2.153_{\pm{0.436}}$      \\[3pt]
\bottomrule[2pt]
\end{tabular}
    \vspace{3mm} \\
    
\centering
\begin{tabular} {ccccccccc }
\toprule[2pt]
        $ \sqrt{\epsilon_{\rm PEHE}} $         & \multicolumn{2}{c}{ \bf{IHDP} 10\%}        &              & \multicolumn{2}{c}{\bf{IHDP} 20\%}& &  \multicolumn{2}{c}{\bf{IHDP} 40\%}      \\[3pt]\cmidrule{2-3}\cmidrule{5-6}\cmidrule{8-9}
Method                      & labeled  & unlabeled & \multicolumn{1}{l}{} & labeled & unlabeled & \multicolumn{1}{l}{} & labeled & unlabeled\\[3pt] \cmidrule{1-9}

CP~($\lambda_o=0$)  &  $^\dagger2.883_{\pm{3.708}}$ & $^\dagger3.004_{\pm{4.071}}$  & \multicolumn{1}{l}{} &  $^\dagger1.972_{\pm{1.930}}$  & $^\dagger2.144_{\pm{2.465}}$  &      &  $1.574_{\pm{1.392}}$  & $1.674_{\pm{1.874}}$      \\[3pt]

CP~($\lambda_e=0$) &        $2.494_{\pm{3.201}}$       &       $2.698_{\pm{3.461}}$         & \multicolumn{1}{l}{} &     $ 1.728_{\pm{2.194}}$   &   $1.977_{\pm{2.450}}$      &      &  $1.344_{\pm{1.383}}$  & $1.585_{\pm{1.923}}$        \\[3pt]

CP &       $\bf 2.427_{\pm{3.189}}$ & $\bf 2.652_{\pm{3.469}}$        & \multicolumn{1}{l}{} &     $\bf 1.686_{\pm{1.838}}$          & $\bf 1.961_{\pm{2.343}}$       &  &   $\bf 1.299_{\pm{1.001}} $& $\bf 1.485_{\pm{1.433}}$      \\[3pt]
\bottomrule[2pt]
    \end{tabular}
        \vspace{3mm} \\
    
% 以下２つ人工データ
\centering
\begin{tabular} {ccccccccc }
\toprule[2pt]
        $ \sqrt{\epsilon_{\rm PEHE}} $         & \multicolumn{2}{c}{ \bf{Synthetic} 10\%}        &              & \multicolumn{2}{c}{\bf{Synthetic} 20\%}& &  \multicolumn{2}{c}{\bf{Synthetic} 40\%}      \\[3pt]\cmidrule{2-3}\cmidrule{5-6}\cmidrule{8-9}
Method                      & labeled  & unlabeled & \multicolumn{1}{l}{} & labeled & unlabeled & \multicolumn{1}{l}{} & labeled & unlabeled\\[3pt] \cmidrule{1-9}

CP~($\lambda_o=0$)  &  $0.334_{\pm{0.112}}$ & $0.367_{\pm{0.124}}$  & \multicolumn{1}{l}{} &  $0.256_{\pm{0.067}}$  & $0.263_{\pm{0.089}}$  &      &  $0.236_{\pm{0.044}}$  & $0.224_{\pm{0.063}}$      \\[3pt]

CP~($\lambda_e=0$) &        $0.320_{\pm{0.110}}$       &       $0.356_{\pm{0.126}}$         & \multicolumn{1}{l}{} &     $ 0.272_{\pm{0.062}}$   &   $0.269_{\pm{0.083}}$      &      &  $0.232_{\pm{0.036}}$  & $0.217_{\pm{0.063}}$        \\[3pt]


CP &       $\bf 0.307_{\pm{0.125}}$ & $\bf 0.331_{\pm{0.149}}$        & \multicolumn{1}{l}{} &     $\bf 0.249_{\pm{0.063}}$          & $\bf 0.259_{\pm{0.086}}$       &  &   $\bf 0.229_{\pm{0.031}} $& $\bf 0.211_{\pm{0.049}}$      \\[3pt]
\bottomrule[2pt]
    \end{tabular}
\end{table}







% ２つのpropagationを比較
Our proposed method has two different propagation terms, the outcome propagation term and the ITE propagation term, as regularizers for semi-supervised learning. 
Table~\ref{table:proposed_comparison} investigates the contributions by the different propagation terms. 
The proposed method using the both propagation terms~(denoted by CP)~shows better results than the one only with the ITE propagation denoted by CP~($\lambda_o=0$); on the other hand, the improvement over the one only with the outcome regularization is marginal.
This observation implies the outcome propagation contributes more to the predictive performance than the ITE propagation.

% hyperparameterへの依存
We also examine the sensitivity of the performance to the regularization hyper-parameters.
% Figure~\ref{fig:hyperparameters} reports the results using $40\%$ and $10\%$ of the whole data as the training data of the News and IHDP datasets, respectively. 
% The proposed method seems rather sensitive to the strength of the regularization terms, particularly on the IHDP dataset, which suggests that the regularization parameters should be carefully tuned using validation datasets in the proposed method. 
%\harada{
Figure~\ref{fig:hyperparameters} reports the results using $10\%$, $20\%$ ,and $40\%$ of the whole data as the training data of the News, IHDP, and Synthetic datasets, respectively. 
The proposed method seems rather sensitive to the strength of the regularization terms, particularly on the IHDP dataset, which suggests that the regularization parameters should be carefully tuned using validation datasets in the proposed method. 
%}
%We also note that the proposed method includes many hyper-parameters which should be fine-tuned. 
In our experimental observations, slight changes in the hyper-parameters sometimes caused significant changes of predictive performance. 
We admit the hyper-parameter sensitivity is one of the current limitations in the proposed method and  efficient tuning of the hyper-parameters should be addressed in future.
% \harada{We compare the proposed method with the state-of-the-art methods varying the magnitude of noise on outcomes. Fig.~\ref{fig:noise_result}~delivers the performance comparison in terms of $\rm \sqrt{\epsilon_{PEHE}}$. Note that the results when $c=1$ correspond to the results in Table~\ref{table:ihdp}~and Table~\ref{table:news}. It is observed that the proposed method is not sensitive when the magnitude of noise is not significant. However, if we add more larger noise, the proposed method fails to propagate outcomes effectively and get more severely affected by the noisy instances than baselines. In both datasets,  we remark that as well as other semi-supervised learning methods, too noisy instances could be crucial as some previous studies have reported~\cite{vahdat2017toward,bui2018neural,du2018robust,liu2012robust}. }

% ノイズを変えて実験
Finally, we compare the proposed method with the state-of-the-art methods by varying the magnitude of noises added to the outcomes. 
Fig~\ref{fig:noise_result}~shows the performance comparison in terms of $\rm \sqrt{\epsilon_{PEHE}}$.
% Note that the results when $c=1$ correspond to the previous results in Tables~\ref{table:news} and \ref{table:ihdp} . 
Note that the results when $c=1$ are the same as those in Tables~\ref{table:news}, \ref{table:ihdp} and \ref{table:synthetic}.
The proposed method stays tolerant of relatively small magnitude of noises; however, with larger label noises, it suffers more from wrongly propagated outcome information than the baselines. 
This is consistent with the previous studies reporting the vulnerability of semi-supervised learning methods against label noises~\cite{vahdat2017toward,bui2018neural,du2018robust,liu2012robust}. 

 \begin{figure}[tb]
%  \begin{minipage}{0.2455\hsize}
\centering
\begin{minipage}{0.43\hsize}
%   \begin{center}
   \includegraphics[width=\linewidth]{images/news_in_reg.pdf}
      \centering
   (a): labeled~(News)
%   \end{center}
 \end{minipage}
 \centering
  \begin{minipage}{0.43\hsize}
%   \begin{center}
%    \includegraphics[width=\linewidth]{images/ihdp_out_reg.pdf}
   \includegraphics[width=\linewidth]{images/news_reg.pdf}
      \centering
   (b): unlabeled~(News)
%   \end{center}
 \end{minipage}\\
 \centering
 \begin{minipage}{0.43\hsize}
%   \begin{center}
   \includegraphics[width=\linewidth]{images/ihdp_in_reg.pdf}
   \centering
   (c): labeled~(IHDP)
%   \end{center}
 \end{minipage}
 \centering
  \begin{minipage}{0.43\hsize}
%   \begin{center}
%    \includegraphics[width=\linewidth]{images/news_out_reg.pdf}
   \includegraphics[width=\linewidth]{images/ihdp_reg.pdf}
      \centering
   (d): unlabeled~(IHDP)
%   \end{center}
 \end{minipage}
  \centering
 \begin{minipage}{0.43\hsize}
%   \begin{center}
   \includegraphics[width=\linewidth]{images/synthetics_in_reg.pdf}
   \centering
   (e): labeled~(Synthetic)
%   \end{center}
 \end{minipage}
 \centering
  \begin{minipage}{0.43\hsize}
%   \begin{center}
%    \includegraphics[width=\linewidth]{images/news_out_reg.pdf}
   \includegraphics[width=\linewidth]{images/synthetics_reg.pdf}
      \centering
   (f): unlabeled~(Synthetic)
%   \end{center}
 \end{minipage}
 \caption{ Sensitivity of the performance to the hyper-parameters.  The colored bars indicate $\sqrt{\epsilon_{\text{PEHE}}}$ for (a)(b) the News dataset, (c)(d) the IHDP dataset  and (e)(f) the Synthetic dataset, when using the largest size of labeled data.
 The deeper-depth colors indicate larger errors.
 It is observed that the proposed method is somewhat sensitive to the choice of the hyper-parameters, especially, the strength of the outcome regularization~($\lambda_o$).  }\label{fig:hyperparameters}
\end{figure}
%
\begin{figure}[t]
 \begin{minipage}{0.33\hsize}
  \begin{center}
   \includegraphics[width=\linewidth]{images/news_noise.pdf}
   (a):  News
  \end{center}
 \end{minipage}
 \begin{minipage}{0.33\hsize}
  \begin{center}
   \includegraphics[width=\linewidth]{images/ihdp_noise.pdf}
   (b):  IHDP
  \end{center}
 \end{minipage}
  \begin{minipage}{0.33\hsize}
  \begin{center}
   \includegraphics[width=\linewidth]{images/synthetic_noise.pdf}
   (c):Synthetic  
  \end{center}
 \end{minipage}
    \caption{Performance comparisons for different levels of noise $c$ added to the labels on (a) News dataset, (b) IHDP dataset, and (c) Synthetic dataset. Note that the results when $c=1$ correspond to the previous results (Tables~\ref{table:news},\ref{table:ihdp}, and \ref{table:synthetic}).}\label{fig:noise_result}
    \end{figure}
    

% レイアウトの問題でここに表を移動
% \begin{table}[tbp]
% \caption{\harada{The performance comparison between the proposed methods on News dataset. $\lambda_o=0$ and $\lambda_e=0$ indicate CP with only treatment effect regularization and outcome regularization. $^\dagger$ indicates that our proposed method (CP) performs statistically significantly better than the baselines by the $t$-test ($p<0.05)$. The bold results indicate the best results in terms of the average.}\label{table:cp_comparison}}
% \centering
% \begin{tabular} {ccccccccc }
% \toprule[2pt]
%         $ \sqrt{\epsilon_{\rm PEHE}} $         & \multicolumn{2}{c}{ \bf{News} 1\%}        &              & \multicolumn{2}{c}{\bf{News} 5\%}& &  \multicolumn{2}{c}{\bf{News} 10\%}      \\[3pt]\cmidrule{2-3}\cmidrule{5-6}\cmidrule{8-9}
% Method                      & labeled  & unlabeled & \multicolumn{1}{l}{} & labeled & unlabeled & \multicolumn{1}{l}{} & labeled & unlabeled\\[3pt] \cmidrule{1-9}

% CP~($\lambda_o=0$)  &  $\bf 2.812_{\pm{651}}$ & $2.806_{\pm{0.598}}$  & \multicolumn{1}{l}{} &  $^\dagger2.527_{\pm{0.474}}$  & $\dagger2.531_{\pm{0.523}}$  &     &  $^\dagger2.400_{\pm{0.347}}$ & $^\dagger2.410_{\pm{0.450}}$  \\[3pt]

% CP~($\lambda_e=0$) &        $2.879_{\pm{0667}}$       &       $2.885_{\pm{0.609}}$         & \multicolumn{1}{l}{} &     $ 2.351_{\pm{0.450}}$   &   $2.483_{\pm{0.481}}$      &       &  $\bf 1.996._{\pm{0.338}}$ & $2.221_{\pm{0.455}}$   \\[3pt]

% CP &        $2.844_{\pm{0.683}}$       &       $\bf 2.823_{\pm{0.656}}$        & \multicolumn{1}{l}{} &     $\bf 2.310_{\pm{0.430}}$          & $\bf 2.446_{\pm{0.471}}$       &  &   $2.003_{\pm{0.393}} $& $\bf 2.153_{\pm{0.436}}$      \\[3pt]
% \bottomrule[2pt]
%     \end{tabular}
% \end{table}




% 実データ→人工データの順番がいいかも?
% 人工と実でそれぞれサブせくしょんにした方が良いか

% \subsection{Experiments on synthetic dataset}
% \harada{Aiming to examine the robustness against the noisy instances, we conduct experiments using synthetic dataset. In this experiments, we vary the magnitude and proportion of noisy outcomes.}

% \subsection{RQ3: Is the proposed model robust to noisy outcomes?}
% \subsubsection{Experimental setting.}

% \harada{
% We generate $1,000$ instances who have $8$ covariates ${\bf x}$ from $\mathcal{N}({\bf 0}^{8\times1},0.5\times(\Sigma+\Sigma^{\top} ))$ where $\Sigma $ $\sim\mathcal{U}((-1,1)^{8\times8})$ and the treatment $t$ from $t\mid\text{Bern}(\sigma({\bf w}_t^{\top}{\bf x}+\epsilon))$ where $\sigma(x)=\frac{1}{1+\exp(-x)}, {\bf w}_t\sim\mathcal{U}((-1,1)^{8\times1}), \epsilon_t\sim\mathcal{N}(0, 0.1)$. We consider the following two types of outcome generative systems:~(i): a non-linear outcome and treatment effect model~(ii): a linear outcome and treatment effect model.
% (i): In the treatment outcome and control outcome is generate as $y^{1}\mid x\sim ({\bf w}^{\top}_{y^{1}}x +c\epsilon_{y})$ and $y^{0}\mid x\sim ({\bf w}^{\top}_{y^{0}}x+c\epsilon_{y})$, where ${\bf w} \sim\mathcal{U}((-1,1)^{8\times1}), \epsilon_y\sim\mathcal{N}({\bf 0}, 0.1)$.  (ii):In the treatment outcome and control outcome is generate as  $y^{1}\mid x\sim (\sin({\bf w}^{\top}_yx) +c\epsilon_{y})$ and $y^{0}\mid x\sim (\cos({\bf w}^{\top}_yx) +c\epsilon_{y})$, where ${\bf w} \sim\mathcal{U}((-1,1)^{8\times1}), \epsilon_y\sim\mathcal{N}({\bf 0}, 0.1)$. $c$ is the parameter which controls the magnitude of noise. 
% Varying the magnitude of noise, we observe how the performance change.
% }
% As we focus on the effect of data scarcity issue, we use $10\%$ of whole data as train data and validation data and use the rest of them as test data.

% \harada{Based on the previous works~\cite{hill2011bayesian,johansson2016learning}, we consider more noisy outcomes for each dataset. In addition to the noiseless potential outcomes, we add noise $c\epsilon$ where $c$ is the magnitude of noise and $\epsilon\sim\mathcal{N}(0, 1)$. Then we vary the magnitude of noise $c$ from $1,3,5,7,9$ and observe how the performances change.
% Note that in both dataset, the outcomes employed in the previous studies~\cite{hill2011bayesian,johansson2016learning,shalit2017estimating} and RQ1  are generated when $c=1$. See the original papers for more detail on the outcome generation~\cite{hill2011bayesian,johansson2016learning}.}

% As we focus on the effect of data scarcity issue, we use $10\%$ of whole data as train data for IHDP dataset and $1\%$ for News dataset.

% \subsubsection{Results and discussions}
% % 人工データ用の説明
% \harada{We compare our method with the state-of-the-art methods varying the magnitude of noisy outcomes. Fig.~\ref{fig:synthetic_result}~demonstrate the performance comparison in terms of $\rm \epsilon_{PEHE}$. It is observed that the proposed method is not sensitive when the magnitude of noise is not significant but if we add more crucial noise, the proposed method seem to fail to propagate outcomes effectively and get more severely affected by the noisy instances. In case~(i), the proposed method yield better performances than baselines when the magnitude is small, however when the magnitude is large~($c=9$), the proposed method is not capable of propagations effectively. In case~(ii), similar phenomena is observed when small magnitude. Standard deviation gets higher than baselines. Though our framework is robust to noisy instances to some extent, we remark that as well as other semi-supervised learning methods, too noisy instances could be more crucial than baselines  some previous studies have noted~\cite{vahdat2017toward,bui2018neural}.
% We do not observe the statistically significant difference between TARNet and CFR.}




% IHDPとNews用の説明
% \harada{We compare the proposed method with the state-of-the-art methods varying the magnitude of noisy outcomes. Fig.~\ref{fig:noise_result}~demonstrate the performance comparison in terms of $\rm \epsilon_{PEHE}$. Note that the results when $c=1$ correspond to the results in Table~\ref{table:ihdp}~and Table~\ref{table:news}. It is observed that the proposed method is not sensitive when the magnitude of noise is not significant. However, if we add more crucial noise, the proposed method seem to fail to propagate outcomes effectively and get more severely affected by the noisy instances than baselines. In both datasets,  we remark that as well as other semi-supervised learning methods, too noisy instances could be more crucial than baselines as some previous studies have reported~\cite{vahdat2017toward,bui2018neural}. }



% It is observed that existing methods which rely only on the labeled data do not seem to work properly . In contrast, the proposed method shows the robustness against the the scarcity of labeled data and effectively. 
% With regard to noisy instances, the proposed method is not sensitive when the magnitude of noise is not significant. However, if we add more crucial noise, the proposed method seem to fail to propagate outcomes effectively and get more severely affected by the noisy instances. In case~(i), the proposed method get degraded given even small proportion of noise~(Fig.~\ref{fig:synthetic_result}(a), (b), (c))~and baselines also failed when large proportion and large noise. We do not observe the statistically significant difference between TARNet and CFR. In case~(ii), the proposed method is not sensitive when the noisy instances are less~(Fig.~\ref{fig:synthetic_result_sin}(a), (b)). However, the proposed method seems to get more affected by noisy instances when the proportion of noise increase and could be crucial if more instances get noisy. CFR-Wass show better results than TARNet. We assume more complex outcome systems of case~(ii) than case~(i)~make such difference.



% \begin{figure}[t]
%  \begin{minipage}{0.33\hsize}
%   \begin{center}
%    \includegraphics[width=\linewidth]{images/synthetic_5.pdf}
%    (a):  $5\%$ of noisy instances
%   \end{center}
%  \end{minipage}
%  \begin{minipage}{0.33\hsize}
%   \begin{center}
%    \includegraphics[width=\linewidth]{images/synthetic_10.pdf}
%    (b):  $10\%$ of noisy instances
%   \end{center}
%  \end{minipage}
%   \begin{minipage}{0.33\hsize}
%   \begin{center}
%    \includegraphics[width=\linewidth]{images/synthetic_20.pdf}
%    (c):  $20\%$ of noisy instances
%   \end{center}
%  \end{minipage}
%    \caption{Performance comparison on synthetic dataset across 10 times experiments in case~(i). Shaded space indicate the standard deviation.}\label{fig:synthetic_result}
%     \begin{minipage}{0.33\hsize}
%   \begin{center}
%    \includegraphics[width=\linewidth]{images/synthetic_sin_5.pdf}
%    (a): $5\%$ of noisy instances
%   \end{center}
% %   \caption{Performance comparison on synthetic dataset across 10 times experiments. Each error bar indicate the standard deviation.}\label{fig:synthetic_result}
%  \end{minipage}
%  \begin{minipage}{0.33\hsize}
%   \begin{center}
%    \includegraphics[width=\linewidth]{images/synthetic_sin_10.pdf}
%    (a): $10\%$ of noisy instances
%   \end{center}
%  \end{minipage}
%   \begin{minipage}{0.33\hsize}
%   \begin{center}
%    \includegraphics[width=\linewidth]{images/synthetic_sin_20.pdf}
%    (a): $20\%$ of noisy instances
%   \end{center}
%  \end{minipage}
%    \caption{Performance comparison on synthetic dataset across 10 times experiments in case~(ii). Shaded space indicate the standard deviation.}\label{fig:synthetic_result_sin}
% \end{figure}


% 人工データ
    
% \begin{figure}[t]
%  \begin{minipage}{0.49\hsize}
%   \begin{center}
%    \includegraphics[width=\linewidth]{images/synthetic_case1.pdf}
%    (a):  Case~(i)
%   \end{center}
%  \end{minipage}
%  \begin{minipage}{0.49\hsize}
%   \begin{center}
%    \includegraphics[width=\linewidth]{images/synthetic_case2.pdf}
%    (b):  Case~(ii)
%   \end{center}
%  \end{minipage}
%    \caption{Performance comparison on synthetic dataset across 10 times experiments in case~(ii). Shaded space indicate the standard deviation.}\label{fig:synthetic_result}
% \end{figure}



% % Please add the following required packages to your document preamble:
% % \usepackage{multirow}
% \begin{table}[]
% \begin{tabular}{cccccc}
%         & \multicolumn{2}{c}{\multirow{2}{*}{\textbf{IHDP 1\%}}} &  & \multicolumn{2}{c}{\multirow{2}{*}{\textbf{IHDP 5\%}}} \\
%         & \multicolumn{2}{c}{}                                   &  & \multicolumn{2}{c}{}                                   \\
% method  & Within-sample             & Without-sample             &  & Within-sample             & Without-sample             \\
% OLS-1   &                           &                            &  &                           &                            \\
% OLS-2   &                           &                            &  &                           &                            \\
% Jaccard &                           &                            &  &                           &                            \\
% katz    &                           &                            &  &                           &                           
% \end{tabular}
% \end{table}
% \clearpage

% \begin{figure}[bt]
%  \begin{minipage}{0.33\hsize}
%   \begin{center}
%    \includegraphics[width=\linewidth]{images/synthetic.pdf}
%   \end{center}
% %   \caption{Performance comparison on synthetic dataset across 10 times experiments. Each error bar indicate the standard deviation.}\label{fig:synthetic_result}
%  \end{minipage}
%  \begin{minipage}{0.33\hsize}
%   \begin{center}
%    \includegraphics[width=\linewidth]{images/synthetic.pdf}
%   \end{center}
%  \end{minipage}
%   \begin{minipage}{0.33\hsize}
%   \begin{center}
%    \includegraphics[width=\linewidth]{images/synthetic.pdf}
%   \end{center}
%  \end{minipage}
%    \caption{Performance comparison on synthetic dataset across 10 times experiments. Each error bar indicate the standard deviation.}\label{fig:synthetic_result}
%     \begin{minipage}{0.33\hsize}
%   \begin{center}
%    \includegraphics[width=\linewidth]{images/synthetic.pdf}
%   \end{center}
% %   \caption{Performance comparison on synthetic dataset across 10 times experiments. Each error bar indicate the standard deviation.}\label{fig:synthetic_result}
%  \end{minipage}
%  \begin{minipage}{0.33\hsize}
%   \begin{center}
%    \includegraphics[width=\linewidth]{images/synthetic.pdf}
%   \end{center}
%  \end{minipage}
%   \begin{minipage}{0.33\hsize}
%   \begin{center}
%    \includegraphics[width=\linewidth]{images/synthetic.pdf}
%   \end{center}
%  \end{minipage}
%    \caption{Performance comparison on synthetic dataset across 10 times experiments. Each error bar indicate the standard deviation.}\label{fig:synthetic_result}
% \end{figure}


% \begin{figure}[t]
%  \begin{minipage}{\hsize}
%   \begin{center}
%    \includegraphics[width=\linewidth,height=200pt]{images/synthetic.pdf}
%   \end{center}
%  \end{minipage}
%    \caption{Performance comparison on synthetic dataset across 10 times experiments. Each error bar indicate the standard deviation.}\label{fig:synthetic_result}
% \end{figure}



\section{Related work\label{sec:related}}
%We review related work in terms of treatment effect estimation and semi-supervised learning.
\subsection{Treatment effect estimation}
Treatment effect estimation has been one of the major interests in causal inference and widely studied in various domains.
Matching~\cite{rubin1973matching,abadie2006large} is one of the most basic and commonly used treatment effect estimation techniques. 
It estimates the counterfactual outcomes using its nearby instances, whose idea is similar to that of graph-based semi-supervised learning. 
Both methods assume that similar instances in terms of covariates have similar outcomes.   
% To mitigate the curse of dimensionality problem in matching, the propensity matching method relying on one-dimensional propensity score was proposed~\cite{rosenbaum1983central}. 
% The propensity score is the probability of an instance to get treatment, and is also used for debiasing the treatment biases~\cite{}.
To mitigate the curse of dimensionality and selection bias in matching, the propensity matching method relying on the one-dimensional propensity score was proposed~\cite{rosenbaum1983central,rosenbaum1985constructing}. 
The propensity score is the probability of an instance to get a  treatment, which is modeled using probabilistic models like logistic regression, and has been successfully applied in various domains to estimate treatment effects unbiasedly~\cite{lunceford2004stratification}. 
% Tree-based methods can build an expressive model based on many weak learners~\cite{chipman2010bart,wager2018estimation}.
Tree-based methods such as regression trees and random forests have also been well studied for this problem~\cite{chipman2010bart,wager2018estimation}.
%and they make many subsets of instances and learn propensity tree. 
%They use regression trees divide the covariate space into 
%The leaves of trained trees are uzed to estimate treatment effect~\cite{wager2018estimation}. 
One of the advantages of such models is that they can build quite expressive and flexible models to estimate treatment effects.
% Recently, deep learning-based methods have been successfully applied for treating the effect estimation~\cite{shalit2017estimating,johansson2016learning}. 
Recently, deep learning-based methods have been successfully applied to treatment effect estimation~\cite{shalit2017estimating,johansson2016learning}. 
Balancing neural networks~(BNNs)~\cite{johansson2016learning} aim to obtain balanced representations of a treatment groups and a control group by minimizing the discrepancy between them, such as the Wasserstein distance~\cite{shalit2017estimating}. 
Most recently, some studies have addressed causal inference problems on network-structured data~\cite{guo2020learning,alvari2019less,veitch2019using}.  
Alvari et al. applied the idea of manifold regularization using users activities as causality-based features to detect harmful users in social media~\cite{alvari2019less}. 
Guo et al. considered treatment effect estimation on social networks using graph convolutional balancing neural networks~\cite{guo2020learning}. 
In contrast with their work assuming the network structures are readily available, we do not assume them and considers matching network defined using covariates.

% knn~\cite{lunceford2004stratification}

\subsection{Semi-supervised learning}
% Semi-supervised learning have been a very popular method. Thus it is impossible to cover all works, we explain 
Semi-supervised learning, which exploits both labeled and unlabeled data, is one of the most popular approaches, especially in scenarios when only limited labeled data can be accessed~\cite{bengio2007greedy,hinton2006fast}.
Semi-supervised learning has many variants, and because it is almost impossible to refer to all of them, we mainly review the graph-based regularization methods, known as label propagation or manifold regularization~\cite{zhu2002learning,belkin2006manifold,weston2012deep}. 
Utilizing a given graph or a graph constructed based on instance proximity, graph-based regularization encourages the neighbor instances to have similar labels or outcomes~\cite{zhu2002learning,belkin2006manifold}. 
%This technique is called {\it label propagation} because the key idea behind this technique is that the neighbors of each instance are regularized to have the same label. 
Such idea is also applied to representation learning in deep neural networks~\cite{weston2012deep,yang2016revisiting,kipf2016semi,iscen2019label,bui2018neural}; they encourage nearby instances not only to have similar outcomes, but also have similar intermediate representations, which results in remarkable improvements from ordinary methods. 
One of the major drawbacks of semi-supervised approaches is that label noises in training data can be quite harmful; therefore, a number of studies managed to mitigate the performance degradation~\cite{vahdat2017toward,pal2018label,du2018robust,liu2012robust}.

One of the most related work to our present study is graph-based semi-supervised prediction under sampling biases of labeled data~\cite{zhou2019graph}.
%The graph-based semi-supervised problem under the sampling biases of labeled data is one of the most related works~\cite{zhou2019graph}.
The important difference between this work and ours is that they do not consider intervention and we do not consider the sampling biases of labeled data. 






\section{Conclusion}
We addressed the semi-supervised ITE estimation problem.
In comparison to the existing ITE estimation methods that only rely on labeled instances including treatment and outcome information, we proposed a novel semi-supervised ITE estimation method that also utilizes unlabeled instances.
The proposed method called counterfactual propagation is built on two ideas from causal inference and semi-supervised learning, namely, matching and label propagation, respectively; accordingly, we devised an efficient learning algorithm.
Experimental results using the semi-simulated real-world datasets revealed that our methods performed better in comparison to several strong baselines when the available labeled instances are limited. 
% However, at the same time, reasonable similarity design and hyper-parameter tuning remained issues.
However, this method had issues related to reasonable similarity design and hyper-parameter tuning.


%Our future work includes addressing the biased distribution of labeled instances. 
%As mentioned in the section discussing the related work
%In this paper, we did not consider such sampling biases for labeled data, but some debiasing techniques~\cite{zhou2019graph} can be successfully integrated into our framework. 

% 順番を変更
One of the possible future directions is to make use of balancing techniques such as the one used in CFR~\cite{shalit2017estimating}, which can be also naturally integrated into our model.
Our future work also includes addressing the biased distribution of labeled instances. 
As mentioned in Related work, we did not consider such sampling biases for labeled data. 
Some debiasing techniques~\cite{zhou2019graph} might also be successfully integrated into our framework. 
In addition, robustness against noisy outcomes under semi-supervised learning framework is still the open problem and will be addressed in the future.



%
% ---- Bibliography ----
%
% BibTeX users should specify bibliography style 'splncs04'.
% References will then be sorted and formatted in the correct style.
%
% \bibliographystyle{splncs04}
% \bibliography{mybibliography}
% \clearpage

\bibliographystyle{splncs04}
\bibliography{MAIN}

\end{document}
